\documentclass[11pt,a4paper]{article}

\usepackage[T1]{fontenc}
\usepackage[utf8]{inputenc}
\usepackage[margin=2.3cm]{geometry}
\usepackage{graphicx}
\usepackage{xcolor}
\usepackage{amsmath}
\usepackage{caption}
\usepackage{parskip}
\usepackage[hidelinks]{hyperref}

\captionsetup{font=small,labelfont=bf,width=0.95\textwidth}
\setlength{\parskip}{0.6em}

\definecolor{boxgrey}{HTML}{F2F5F8}
\definecolor{inkblue}{HTML}{12263A}

\title{\vspace{-1.2cm}\textbf{Catch-up report: power side-channel input recovery\\on the CW305 board}}
\author{Tarek Alhafez}
\date{30 August 2026\\
\small(follow-up work of 31 August in Section 7, of 1 September in Section 8)}

\begin{document}
\maketitle
\vspace{-0.8cm}

\section*{Why I am writing this}

This report is a catch-up for my professor. The project was stuck for a long time, and
the older documents in the repository do not explain the real reason. I want to put in
one place: what the project does, what was actually broken, what I fixed, and what the
new measurements say. I also ran a set of harder tests today, because the earlier
results were measured only in easy conditions and I did not trust them enough.

I write it in a normal way, not in the style of a paper, so it is faster to read.

\section*{The short version}

\begin{itemize}
\item The attack now works on real hardware. The power traces come from a Xilinx
      Artix-7 FPGA on a CW305 board, captured with a ChipWhisperer-Lite. Nothing in
      the measurement is simulated.
\item The old problem was not the algorithm. The FPGA was still running the
      \emph{stock AES example} instead of my design. So the host was writing MNIST
      pixels into AES registers, and the traces were AES power. This is why detection
      looked low and inconsistent for weeks.
\item After fixing this, the same reconstruction code recovers the input. With the
      search radius tuned (Section 8.1) the foreground F1 is $0.969$ and all 40 test
      images are recognised correctly. The $0.843$ quoted through Sections 6 and 7 was
      measured at the search radius carried over from the paper, which turned out to
      be badly wrong for this hardware.
\item Averaging matters a lot. With one hardware repeat the attack is close to the
      trivial baseline. The signal of the convolution unit is small compared to the
      rest of the chip, so repeats are needed to bring it out.
\item The template does not survive every change of condition, and the pattern is more
      specific than I expected. It survives a change of amplifier gain in both
      directions, but it breaks when the clock frequency is changed after profiling, and
      at 10\,MHz it fails completely. The reprogramming row of Table~2 has to be
      withdrawn; Section 7.2 explains why.
\item The attacker cannot profile with any probe kernels he likes. When I profiled with
      one-hot kernels and attacked traces from the real trained kernels, the attack
      became worse than answering ``everything is background''.
\item Going from 40 to 300 profiling images made the attack \emph{worse} at a fixed
      search radius, which turned out to be my own tuning mistake rather than a
      property of the attack. Tuned, the larger template wins: $0.969$ against
      $0.958$. The best radius shrinks as the template grows, so the two have to be
      tuned together.
\item A template is bound to the clock it was profiled at, and the binding is
      asymmetric: in all six pairs of distinct clocks, profiling above and attacking
      below beats the reverse, and it holds only inside a band of roughly a factor of
      two. Section 8.3 replaces the ``profile at the highest clock'' rule I wrote from
      two frequencies, which a four-frequency grid showed to be wrong. For the 5 and
      10\,MHz pair, Section 8.7 then shows the asymmetry is not caused by the template
      match at all --- the failing direction retrieves and ranks the true patch
      \emph{better} than the working one, at every radius tested, so the loss happens
      in the reconstruction stage.
\item Randomising the accelerator clock is a much weaker defence than I proposed. It
      costs an attacker who profiled at 5\,MHz half his F1, but an attacker who
      profiled at 10\,MHz keeps $0.737$ of $0.882$ and still reads three digits in four,
      because that template sits within transfer range of the whole band (Section 8.4).
\item Mixing conditions in one template does not help at 10\,MHz even when the coverage
      is matched, though at 5\,MHz it ties; profiling at a mid-range clock with a tuned
      radius beats it at both clocks using half the data (Section 8.5).
\item A template does not transfer to a different dataset, and I can say why. An
      MNIST-built template attacking Fashion-MNIST scores at chance ($0.034$ above a
      label-blind predictor) while a Fashion-built template on the same traces reaches
      $0.662$. The template is a dictionary of observed $3\times3$ patches --- MNIST
      covers 156 of 512, Fashion 450, with MNIST very nearly a subset --- so transfer
      works only in the direction of increasing coverage (Section 9).
\item Fashion-MNIST is recoverable from power, which no template attack had shown
      before. In greyscale, a classifier reads the recovered clothing at $0.827$ against
      $0.830$ on the originals, so the reconstruction is nearly lossless in
      classification terms (Section 9.4).
\item The attack does not need the leakage amplifier at all. With the bank removed
      entirely --- one flop registering the output, against $72$ in the smallest amplified
      build and $4536$ in the largest --- it reaches MSSIM $0.724$ and recognition $0.746$
      on the paper's $60\,000$-trace profiling set, within noise of the $72$-flop build's
      $0.725$ and $0.757$. What the amplifier changes is how much profiling data is
      needed, not whether the attack is possible: shrinking it $63$ times costs about what
      shrinking the profiling set $37$ times costs, and either loss can be partly bought
      back with the other. For a defender that is the worse reading, because profiling
      traces are the cheap resource here (Sections 9.5 to 9.7).
\item Reducing the victim's window dwell from $8$ cycles to $2$, so it streams closer to
      the way a real accelerator would, did not weaken the attack. Recognition rose from
      $0.746$ to $0.807$ and MSSIM from $0.724$ to $0.762$, on the same test images with
      the amplifier absent in both cases --- about six standard errors, not noise. I
      predicted the opposite before running it. The two bench artifacts I introduced and
      have now removed, the leakage amplifier and the slow dwell, were both things the
      attack did not need (Section 9.8).
\item Three claims in earlier versions of this report were wrong about that, and all three
      failed in the same direction --- treating a limitation of my own build as a fact
      about the hardware. A working note said the attack had no signal without the
      amplifier. The first reading of the sweep measured every width at one small
      profiling budget and concluded the amplifier was close to a precondition. And
      Table \ref{tab:papersettings} recorded the amplifier as ``$72$ flops minimum'',
      where ``minimum'' described a build of mine that would not come up. The cause of
      that dead build was a combinational path from the convolution datapath to a
      top-level pad, which is mine and not the board's (Section 9.7).
\item Run with Power2Picture's own training protocol --- their generator architecture,
      optimiser, loss, batch size and $60\,000$-image profiling set --- the attack reaches
      MSSIM $0.861$ and recognition $0.919$ against a ceiling of $0.898$. That is well
      above their published Fashion-MNIST figures, but it is not a fair comparison and I
      do not present it as one: they attack a full LeNet remotely through on-chip sensors,
      this attacks a single convolution layer through an external shunt (Section 9.6).
\item Raw F1 is not comparable between datasets: the label-blind ceiling moves from
      $0.204$ on MNIST to $0.474$ on Fashion. Every cross-dataset number in this report
      is reported against its own chance baseline for that reason.
\item The generative (Power2Picture) route also works on the corrected captures, with a
      foreground F1 of about $0.82$. The earlier generative number in the repository was
      measured before the authenticity gate existed and must not be quoted.
\item Removing the input binarisation lets real grey levels come back from power alone:
      mean absolute error under $10$ on the 0--255 scale and correlation $0.93$ with the
      true image, from a design confirmed bit-exact on hardware. But a control bitstream
      (Section 8.6) shows I first credited this to the wrong cause: most of the apparent
      gain came from the grey design's wider leakage amplifier, not from the input width.
      With the amplifier held equal, grey is \emph{worse} than binary on foreground F1
      and better only on the intensity measures, which is the narrower claim that
      survives.
\end{itemize}

\section*{1. What the project is doing}

I am reproducing the attack of Wei et al., \emph{I Know What You See} (ACSAC 2018). The
idea of that paper is that the first convolution layer of a binarised neural network
accelerator processes the input image pixel by pixel. Because the power consumption of
the chip depends on the pixels that are inside the convolution window at that moment,
an attacker who measures the power can recover the input image. The attacker never
sees the image, only the current of the chip.

The original paper used a SAKURA-G board with a Spartan-6 FPGA and a 2.5\,GS/s
oscilloscope. I do not have this equipment. I use a CW305 board with an Artix-7
XC7A100T, and a ChipWhisperer-Lite for the capture. Figure~\ref{fig:chain} shows the
measurement chain that I use.

\begin{figure}[h]
  \centering
  \includegraphics[width=\linewidth]{figures/fig01_signal_chain.png}
  \caption{The measurement chain. The convolution runs on the FPGA. The PC only writes
  the input, starts the operation, and does the analysis afterwards. The power signal
  never passes through the PC before the ADC.}
  \label{fig:chain}
\end{figure}

One point is important here, because it was the main doubt in the previous meetings.
The convolution is really computed on the FPGA. The host does not compute it. Before
every capture, the script writes one image and one kernel to the board, reads back all
676 output values, and compares them with a CPU model. If even one value is different,
the capture is refused. In every capture used in this report the result was
$676/676$ bit-exact with zero mismatches.

\section*{2. What was actually blocking the project}

This is the part I want to explain clearly, because the old reports in the repository
give the wrong reason.

On 29 August I found that the board was answering as the stock AES design. The CW305
register pages 2, 3 and 4 returned \texttt{0x02}, \texttt{0x05}, \texttt{0x2e}, which is
the signature of the AES example that ships with the ChipWhisperer. The custom
bitstream was built correctly and met timing, but it was not the image running on the
FPGA.

The consequence is serious. On the AES design those same register pages mean clock
settings, an LED register, a GO register and the AES text input. So when the host wrote
a MNIST image and a convolution kernel, it was writing into AES registers, and the
captured trace was the power of an AES encryption. The reconstruction code was then
trying to find image structure in a trace that contained no image at all. This explains
the weak and unstable detection we saw before.

I added a check that reads those register pages before any capture and refuses to run
if the AES signature is there. Figure~\ref{fig:gate} shows the two days side by side.
I think this is the strongest evidence I have, because the check is not something that
always passes: it failed on 29 August and it passed on 30 August.

\begin{figure}[h]
  \centering
  \includegraphics[width=\linewidth]{figures/fig02_authenticity_gate.png}
  \caption{The same gate on two days. A test that never fails proves nothing. This one
  refused to capture on 29 August and allowed the capture on 30 August.}
  \label{fig:gate}
\end{figure}

There are three different reasons the earlier work looked weak, and I want to keep them
separate, because they need different solutions. Figure~\ref{fig:stalled} summarises
them.

\begin{figure}[h]
  \centering
  \includegraphics[width=\linewidth]{figures/fig10_why_stalled.png}
  \caption{Three separate causes. Only the first one is now solved. The second one is a
  physical limit of the board and is the reason averaging is needed. The third one is a
  documentation problem: the old file \texttt{FINAL\_RESULTS.md} blames per-cycle
  extraction and smearing, and that was not the cause.}
  \label{fig:stalled}
\end{figure}

I want to be direct about the third point. The file \texttt{FINAL\_RESULTS.md} and
\texttt{REPORT\_1} in the repository are older than the corrections, and they explain
the failure in a way that I now believe is wrong. If somebody reads them first, they
will optimise the wrong thing. The current explanation is the one in this report.

\section*{3. What the signal looks like}

Figure~\ref{fig:raw} shows a real captured trace. The FPGA runs at 5\,MHz and the ADC
of the ChipWhisperer samples at four samples per FPGA clock, locked to the same clock.
Each convolution window takes 8 clock cycles, so it takes 32 ADC samples. There are
676 windows, which gives 21632 useful samples out of the 22144 that are captured.

\begin{figure}[h]
  \centering
  \includegraphics[width=\linewidth]{figures/fig03_raw_trace.png}
  \caption{Real captured power signal for one probe kernel and one image. Top: all 676
  windows, drawn as the envelope of the 32 samples inside each window. Bottom left: the
  raw samples of four windows, where the window boundaries are visible. Bottom right:
  the single number that the attack keeps for each window.}
  \label{fig:raw}
\end{figure}

Here is one advantage of my setup compared to the paper. The paper needed a whole
section (their Section 5) for DC restoration and per-cycle curve fitting, because their
oscilloscope sampled asynchronously and the per-cycle power was smeared. My ADC is
phase-locked to the FPGA clock, so I read the per-cycle power directly. I did not
reproduce their Section 5, and I do not need to. This is a simplification, not a
shortcut, and it does not remove any information the attack uses.

From the trace I compute one feature per convolution window: the mean absolute value of
the 32 samples. With 9 probe kernels this gives a vector of 9 numbers for every window.
Figure~\ref{fig:feat} shows that the digit is already visible in this feature, before
any attack algorithm runs. For me this was the moment where I became confident that the
leakage is real.

\begin{figure}[h]
  \centering
  \includegraphics[width=\linewidth]{figures/fig04_feature_map.png}
  \caption{From the captured signal to the attack feature. The FPGA receives a
  binarised image, not the grey image. The power feature map of a single kernel already
  shows the shape of the digit.}
  \label{fig:feat}
\end{figure}

\section*{4. How the attack works in my implementation}

The attack has two stages, and Figure~\ref{fig:algo} shows both.

In the profiling stage I use images that I know. For every convolution window I store
the 9-kernel power vector together with the true $3\times3$ binary patch that produced
it. There are $2^9 = 512$ possible patches. With 40 profiling images I get
$40 \times 676 = 27040$ rows in the template.

In the attack stage I only have traces. The ground truth image is not present in the
attack files at all: the code refuses to run if the attack bundle contains an
\texttt{image} or \texttt{label} array. For every window I take the measured power
vector, split the 9 kernels into 3 groups of 3, search each group in the template with
a kD-tree within a distance $\delta = 1.0$, and intersect the three results. This gives
a list of candidate patches for that window. Then a greedy selection, which follows
Algorithm 2 of the paper, chooses one candidate per window by preferring the candidate
that agrees with its neighbours. Finally the overlapping patches vote on each pixel.

\begin{figure}[h]
  \centering
  \includegraphics[width=\linewidth]{figures/fig05_algorithm.png}
  \caption{The active power-template attack as I implemented it. The grouping,
  $\delta$ and the number of probe kernels are the same as in the paper.}
  \label{fig:algo}
\end{figure}

\section*{5. Where my implementation is different from the paper}

I do not want to claim that I reproduced the paper exactly, because I did not. Some
differences come from the hardware I have, and one difference is important for how the
results should be read. Figure~\ref{fig:geom} and Figure~\ref{fig:cmp} show them.

\begin{figure}[h]
  \centering
  \includegraphics[width=\linewidth]{figures/fig06_geometry.png}
  \caption{The geometry of my layer is not the same as the paper's layer, and my layer
  binarises the pixel before the convolution.}
  \label{fig:geom}
\end{figure}

The most important difference is this one. In the paper the first convolution layer
takes real pixels from 0 to 255 and multiplies them with binary weights. So the power
depends on the grey level, and the attacker can recover a grey image. In my hardware
design the pixel is thresholded at 127 before the convolution, so the layer works on
binary pixels and binary weights. This means that the best I can possibly recover is a
one-bit silhouette of the digit, not the grey levels.

Because of this I score the results against the binarised image, and I report F1, IoU
and MSSIM rather than only pixel accuracy. I also always print the accuracy of the
trivial ``predict everything is background'' answer next to my pixel accuracy. Most
MNIST pixels are background, so pixel accuracy alone looks good even when the attack
recovers nothing. I think this is the honest way to present it.

The second difference is the geometry. The paper uses a $28\times28$ output with
``same'' padding, which is 784 cycles. My hardware does a valid convolution, so the
output is $26\times26$, which is 676 windows. The mapping from window to pixel window,
which is what the attack depends on, is the same idea in both cases.

\begin{figure}[h]
  \centering
  \includegraphics[width=\linewidth]{figures/fig12_comparison.png}
  \caption{My setup next to the setup of the paper.}
  \label{fig:cmp}
\end{figure}

\section*{6. The new and harder tests}

The earlier results were all measured in one comfortable condition: profile and attack
in the same session, same clock, same gain, and with five hardware repeats averaged.
That is not a strong test. Today I ran experiments that are harder, because I want to
know where the attack breaks and not only that it works once.

\subsection*{6.1 How many hardware repeats are really needed}

The old comparison between ``avg 5'' and ``avg 1'' was not a fair comparison, because
those two captures used different images, a different number of profiling images and a
different evaluation set. Three things changed at the same time, so it was not possible
to say that the averaging caused the difference.

I did it properly this time. The capture script stores the individual repeats without
averaging them. So I captured once with five repeats, and then I made the sweep by
using only the first $R$ repeats afterwards. This means every point of the sweep uses
exactly the same images, the same board temperature and the same gain. Only the number
of averaged repeats changes.

\begin{figure}[h]
  \centering
  \includegraphics[width=\linewidth]{figures/fig11_noise_vs_repeats.png}
  \caption{The same measurement, averaged over more repeats. This is the raw power
  feature map, before any attack. The digit appears as noise is averaged away.}
  \label{fig:noiserep}
\end{figure}

\begin{figure}[h]
  \centering
  \includegraphics[width=\linewidth]{figures/fig07_repeat_sweep.png}
  \caption{The averaging sweep. The left plot shows why the baseline must be drawn:
  pixel accuracy is already high at one repeat, but it is close to the trivial answer.
  The metrics that really move are F1 and MSSIM.}
  \label{fig:sweep}
\end{figure}

\begin{table}[h]
  \centering
  \small
  \begin{tabular}{lrrrrrrrr}
\hline
Run & Pixel acc. & All-background & F1 & IoU & MSSIM & Pixel dist. & Recog. & $n$ \\
\hline
1 repeat & 0.919 & 0.888 & 0.494 & 0.328 & 0.501 & 20.6 & 0.65 & 40 \\
2 repeats & 0.938 & 0.888 & 0.644 & 0.475 & 0.623 & 15.7 & 0.70 & 40 \\
3 repeats & 0.946 & 0.888 & 0.698 & 0.536 & 0.673 & 13.8 & 0.78 & 40 \\
4 repeats & 0.958 & 0.888 & 0.778 & 0.637 & 0.752 & 10.8 & 0.93 & 40 \\
5 repeats & 0.960 & 0.888 & 0.792 & 0.656 & 0.769 & 10.2 & 0.90 & 40 \\
\hline
\end{tabular}

  \caption{The averaging sweep, from one single capture subset afterwards. Pixel
  distance is on the 0--255 scale of the paper. ``All-background'' is the pixel accuracy
  of answering that every pixel is background.}
\end{table}

\subsection*{6.2 Does the template survive a change of conditions}

This is the test I consider the most useful, and it is also the closest to a real
attack. In a real attack the profiling and the attack do not happen in the same minute
and under the same settings.

I used one template, built hours earlier from 40 images at 5\,MHz and 40\,dB gain. Then
I captured the \emph{same 40 new images} under six different conditions, so that only
the condition changes and the images do not. The conditions are: the same settings
again as a control, after reprogramming the FPGA completely, with the clock raised to
7.5\,MHz and to 10\,MHz, and with the amplifier gain changed to 30\,dB and to 50\,dB.

Every one of these captures passed the functional check with $676/676$ bit-exact
outputs, which is already a useful result on its own: the datapath is correct even at
10\,MHz.

I did not want to only assume that the six conditions really use the same images, so I
checked it: I took the SHA-256 of the stacked ground-truth images of every evaluation
set, and all of them give the same hash, \texttt{32592af58b9ac9be...}. The
paper-scale evaluation of Section 6.4 uses the same set too. So when two rows of the
tables differ, the images are not the reason.

One question comes immediately here: how can I compare a trace taken at 5\,MHz with a
trace taken at 10\,MHz? The answer is that the ADC clock is derived from the FPGA clock
(\texttt{adc\_src = clkgen\_x4}), so it scales together with it. Every capture has
exactly 22144 samples and 32 samples per convolution window, at every clock. The traces
are therefore directly comparable, and no resampling is done.

\begin{figure}[h]
  \centering
  \includegraphics[width=\linewidth]{figures/fig08_conditions.png}
  \caption{Template transfer under changed conditions. One profile, six capture
  conditions, the same 40 images.}
  \label{fig:cond}
\end{figure}

\begin{table}[h]
  \centering
  \small
  \resizebox{\textwidth}{!}{\begin{tabular}{lrrrrrrrr}
\hline
Run & Pixel acc. & All-background & F1 & IoU & MSSIM & Pixel dist. & Recog. & $n$ \\
\hline
Same conditions (control) & 0.968 & 0.887 & 0.843 & 0.728 & 0.816 & 8.2 & 1.00 & 40 \\
FPGA reprogrammed & 0.964 & 0.887 & 0.818 & 0.693 & 0.790 & 9.3 & 0.97 & 40 \\
Clock 5 $\to$ 7.5\,MHz & 0.903 & 0.887 & 0.273 & 0.158 & 0.382 & 24.7 & 0.38 & 40 \\
Clock 5 $\to$ 10\,MHz & 0.887 & 0.887 & 0.007 & 0.004 & 0.289 & 28.9 & 0.12 & 40 \\
Gain 40 $\to$ 30\,dB & 0.966 & 0.887 & 0.832 & 0.712 & 0.805 & 8.6 & 0.90 & 40 \\
Gain 40 $\to$ 50\,dB & 0.966 & 0.887 & 0.831 & 0.711 & 0.802 & 8.6 & 0.93 & 40 \\
\hline
\end{tabular}
}
  \caption{Transfer of one template to captures made under different conditions. The
  ``FPGA reprogrammed'' row is withdrawn: see Section 7.2. Every other row stands.}
\end{table}

\begin{figure}[h]
  \centering
  \includegraphics[width=0.72\linewidth]{figures/fig09_examples.png}
  \caption{The same six images, recovered under different conditions. The empty row for
  the 10\,MHz capture is not a drawing mistake: at that clock the reconstruction really
  returns an all-background image for every one of the 40 test images.}
  \label{fig:ex}
\end{figure}

The result is easy to read, with one row that I have to withdraw. The row labelled
``FPGA reprogrammed'' gives F1 $0.818$, and I first read that as the template surviving
a reconfiguration. On 31 August I found that reprogramming a CW305 that is already
configured does nothing at all, so that capture almost certainly ran on the same
configuration as the control. The number is real, the label is wrong. Section 7.2
explains what I found and how to test the claim properly.

Changing the clock is a different story. At 7.5\,MHz the F1 falls to about $0.27$, and
at 10\,MHz the attack collapses completely: the reconstruction returns an all-background
image for every test image, so the pixel accuracy is exactly equal to the trivial
baseline and the F1 is almost zero. This is the clearest example in this report of why
the baseline column must be printed. A reader who sees only ``pixel accuracy $0.887$''
would think the attack still works, when in reality it produced nothing at all.

The amplifier gain, on the other hand, turned out not to matter. I changed it in both
directions, down to 30\,dB and up to 50\,dB, and the F1 stayed at $0.832$ and $0.831$
against $0.843$ for the control. I expected some loss here and there is almost none. The
reason is that the feature extraction normalises every trace by its own median and MAD
before the template is queried, so a change of gain is mostly a change of scale and the
normalisation removes it. A change of clock is not a change of scale: it changes the
shape of the power inside the window, and no normalisation can undo that.

So the conclusion is more specific than ``the template is fragile''. The template
tolerates reprogramming and it tolerates a gain change, but it is tied to the clock
frequency it was profiled at.

\subsection*{6.3 Does the attacker need to know the deployed kernels}

In the profiling stage the attacker chooses probe kernels. In my earlier runs the probe
kernels were the same trained kernels that the FPGA was running. This is a strong
assumption, so I wanted to test it. I built a template from captures that used one-hot
probe kernels, and I used it to attack traces that were captured while the FPGA was
running the real trained kernels.

The attack fails completely. The foreground F1 falls to about $0.12$, and the pixel
accuracy falls \emph{below} the all-background baseline, $0.811$ against $0.888$. This
is worse than doing nothing: the reconstruction is actively painting foreground pixels
in places where there is only background. The digit recogniser then reads the recovered
images correctly only one time out of five.

Somebody can object that the one-hot captures and the trained captures were also made at
different times, so maybe the delay caused this and not the probe kernels. I do not
think so, and the control experiment answers it: in Section 6.2 a template was
transferred across a comparable time gap, to images that it had never seen, and it still
reached F1 $0.84$. So the collapse here comes from the probe mismatch, not from the
delay. The practical meaning is that the attacker cannot profile with arbitrary probe
kernels; the profiling must be done with the kernels that are really deployed.

\subsection*{6.4 Paper scale}

All the numbers above use 40 profiling images. The paper uses 300 profiling images and
200 evaluation images. Since the capture turned out to be fast, roughly 0.6 seconds per
image, I also captured 500 images and built a template from 300 of them, which is the
same template size as the paper.

The capture is fast but the attack is not. With 300 profiling images the template has
$300 \times 676 = 202800$ rows instead of $27040$, and the candidate search becomes much
slower for every window.

I also changed the way I evaluate this run, and I think the new way is better science
than what I planned first. Instead of evaluating on 200 new images, I attacked
\emph{the same 40 control images} that were used in Section 6.2. The reason is that a
comparison is only meaningful when one thing changes. If I use a bigger template and
different evaluation images at the same time, I cannot say which of the two caused the
difference. Because this run and the control run share their evaluation set exactly, the
only difference between this row and the control row of Table~\ref{tab:headline} is the
size of the template: 300 profiling images instead of 40.

The result is the opposite of what I expected, and it is the most interesting thing I
found today. \textbf{More profiling data made the attack worse.} With 300 profiling
images the F1 is $0.759$, against $0.843$ with only 40, and the recogniser drops from
40 out of 40 to 37 out of 40. Everything else was identical.

I did not want to only guess the reason, so I looked at what the attack itself records.
Each attack writes the average number of candidate patches per window. With the small
template it is $14.1$ candidates per window; with the large template it is $32.0$, more
than double. This explains the loss. The candidate search keeps every template row that
falls inside the fixed distance $\delta = 1.0$. When the template grows by $7.5$ times,
many more rows fall inside that same radius, so each window returns a longer list of
candidates and the greedy selection has more chances to choose a wrong one.

There is a second number that agrees with this. The ``fallback'' counter, which counts
the windows where nothing at all was found inside $\delta$ and the search had to fall
back to nearest neighbours, drops from $1.00$ to $0.17$ windows per image. So the big
template really does cover the space better. The problem is not coverage, it is that the
threshold was not moved when the template grew.

My conclusion is that $\delta = 1.0$ is not a constant of the attack, it is a value that
has to be tuned together with the size of the template. The paper uses $\delta = 1.0$
with 300 profiling images, but on their hardware the power vector has a different scale,
so their $\delta$ is not directly my $\delta$. I should have tuned it and I did not.
This is a concrete thing to fix next, and it is cheap to test because the captures
already exist and only the attack has to be run again.


\section*{7. Follow-up work, 31 August}

Sections 1 to 6 are the state of 30 August. The three items below were the first three
things on my own next-steps list, and I did them the day after.

\subsection*{7.1 Making the template tolerant to a change of clock}

Section 6.2 showed the template dies when the clock changes. My idea was to profile
under several conditions at once. Before doing that I checked the other idea on the
list, normalising the features per capture, and found it was \emph{already} in the
pipeline: every trace is normalised by its own median and MAD before the template is
queried. That is exactly why a gain change costs nothing, and it is also why it cannot
rescue a clock change: a gain change is a change of scale, which the normalisation
removes, while a clock change alters the shape of the power inside the window.

So I tested the mixing idea properly. I captured 40 fresh profiling images
(5300--5339, disjoint from the evaluation images) twice, once at 5\,MHz and once at
10\,MHz, and built three templates:

\begin{itemize}
\item 5\,MHz only, 40 images;
\item 10\,MHz only, 40 images;
\item mixed, the first 20 images at 5\,MHz plus the \emph{other} 20 at 10\,MHz.
\end{itemize}

All three hold 40 profiling images and the same number of rows. That matters, because
Section 6.4 showed a larger template at a fixed $\delta$ scores worse; if the mixed
template were simply bigger, I could not tell mixing from that effect. Each template
then attacked the same 40 evaluation images at both clocks.

\begin{figure}[h]
  \centering
  \includegraphics[width=\linewidth]{figures/fig13_clock_mixing.png}
  \caption{Three templates, two attack clocks, one evaluation set. Mixing repairs the
  cross-clock case, but profiling only at the higher clock is better still.}
  \label{fig:mixing}
\end{figure}

\begin{table}[h]
  \centering
  \small
  \resizebox{\textwidth}{!}{\begin{tabular}{lrrrrrr}
\hline
Template (40 profile images) & \multicolumn{3}{c}{Attack at 5\,MHz} & \multicolumn{3}{c}{Attack at 10\,MHz} \\
 & F1 & MSSIM & Recog. & F1 & MSSIM & Recog. \\
\hline
Profiled at 5\,MHz only & 0.781 & 0.749 & 0.90 & 0.081 & 0.300 & 0.12 \\
Profiled at 10\,MHz only & 0.803 & 0.639 & 1.00 & 0.560 & 0.484 & 0.80 \\
Mixed, 20 images at each clock & 0.790 & 0.704 & 0.95 & 0.453 & 0.401 & 0.50 \\
\hline
\end{tabular}
}
  \caption{Clock mixing. Every template uses 40 profiling images and the same 40
  evaluation images, so only the profiling clock changes.}
  \label{tab:mixed}
\end{table}

The answer is yes, mixing works: at 10\,MHz the F1 goes from $0.081$ with the 5\,MHz
template to $0.453$ with the mixed one. But the more useful result is the one I did not
expect. \textbf{The template profiled only at 10\,MHz is the best template at both
clocks}: $0.560$ at 10\,MHz, and $0.803$ at 5\,MHz with all 40 images recognised
correctly, which even beats the 5\,MHz template on its own clock.

The failure is not symmetric. Profiling low and attacking high collapses; profiling high
and attacking low transfers almost for free. If this holds at other frequencies, the
practical rule is simple: profile at the highest clock you expect to attack, and do not
bother mixing.

I also looked at why mixing loses to the pure 10\,MHz template. It is not coverage. At
10\,MHz both return the same number of candidate patches per window, $46.5$ against
$46.6$, and both fall back on the nearest neighbour in only $0.03$ windows per image. So
the mixed template finds just as many matches; the 5\,MHz half simply contributes
\emph{wrong} candidates that still fall inside $\delta$. That is the same mechanism as
the paper-scale result of Section 6.4.

One honest limit: each cell is 40 images, so the small gaps, around $0.02$ to $0.05$,
are not resolved. The gap between $0.081$, $0.453$ and $0.560$ clearly is.

\subsection*{7.2 Why the reprogramming row of Table 2 is withdrawn}

While preparing the next experiment I tried to load a different bitstream and the board
kept running the old one. Reprogramming a CW305 whose FPGA is already configured does
nothing: \texttt{cw.target(force=True)}, a direct \texttt{FPGAProgram} call with
\texttt{exceptOnDoneFailure=True}, and even \texttt{eraseFPGA} all report success and
leave DONE high while the old design keeps answering. Only cutting the power blanks the
FPGA; after that, programming works immediately.

That is why the ``FPGA reprogrammed'' capture of Section 6.2 has to be withdrawn. It
ran the capture script without \texttt{--no-program}, so I recorded it as a
reconfiguration, but the FPGA was already configured and almost certainly kept the same
image. The measurement is a valid ``same configuration, some minutes later'' point, not
a reconfiguration test. Testing the real claim needs a power cycle between profiling and
attack, which I will do next.

The authenticity gate does not catch this, and it is worth saying why. The gate compares
register pages 2, 3 and 4 against the stock AES signature. Both of my designs read
\texttt{0x00} there, so the gate can tell my design from the AES example but cannot tell
one of my designs from another. What does distinguish them is writing a byte ramp into
the image region and reading it back: the binary design stores one bit per pixel and
returns $1,0,1,0,\ldots$, while the grey design returns the bytes themselves.

\subsection*{7.3 Removing the binarisation: a grey-pixel design}

Section 5 explains that my hardware thresholds the pixel before the convolution, so the
best I can recover is a silhouette. I therefore built a second design,
\texttt{cw305\_leakage\_grey\_top}, in which the first layer multiplies real 0--255
pixels by binary weights, as in the paper. It is a separate design and the binary one is
untouched, because every number in Sections 1--6 came from the binary design and the two
are not interchangeable: the grey design needs its own profiling from scratch.

Three things were worth learning here.

The host protocol did not have to change. The host already sent 784 bytes, one per
pixel, and the binary register file simply used bit 0 of each. Only the output had to
grow, from 8-bit to 16-bit signed, because a grey MAC reaches $\pm 9\times255 = \pm2295$.
I widened it rather than scaling the result down, since throwing away low bits of the
value I am trying to recover would defeat the point.

The leakage amplifier had to keep riding a kernel-dependent term. The binary core
toggles banks of $\lnot(\text{window} \oplus \text{kernel})$; the grey core toggles the
byte-wise version of the same expression, which is the identical function when a pixel
is one bit. Riding the 16-bit sum instead would have collapsed the per-tap grey
information that the template attack needs.

The first build was wrong in a way that is easy to miss. It reported zero errors and
produced a bitstream, but it did not meet timing: worst slack $-3.142$\,ns with 8440
failing endpoints. Holding the image as one packed 6272-bit vector turned every USB
write into a 784-way enable and every readback into a 784-to-1 byte multiplexer. Storing
the image as a RAM array instead, and having the core ask for the nine window bytes,
fixed it: worst slack $+1.620$\,ns, no failing endpoints, and the LUT count fell from
42\,\% to 25\,\% of the device. I also found that the project constraint file names a
clock net that does not exist, so the convolution had been timed against a 100\,MHz clock
the design never runs at; the grey build names the real net.

Because the board was not available while I did this, I first verified the design in
simulation: a self-checking testbench drives a pseudo-random grey image and compares all
676 outputs against an independent reference, and it passes bit-exact. The host-side
reference model was separately checked against a second implementation.

After the power cycle described in Section 7.2 the grey bitstream loaded, and the design
passes on the real board as well: writing a byte ramp into the image region reads back
all 784 bytes unchanged, which the binary design cannot do, and the functional check
gives $676/676$ outputs bit-exact with the hardware range $-1408 \ldots 643$ equal to
the reference range. So the grey convolution is confirmed on silicon and not only in
simulation.

\subsubsection*{Do the grey levels really come back?}

This is the question that motivated the whole change. The template attack of Section 4
cannot answer it, because it classifies each window into one of $2^9$ binary patches; it
has no way to express a grey value. The generative route can, so I captured 2000 images
with the grey design and trained the same generator as Section 7.4, with the same split
and the same settings as the binary run, changing only the design under attack.

\begin{figure}[h]
  \centering
  \includegraphics[width=\linewidth]{figures/fig14_grey_recovery.png}
  \caption{Recovered images from the two designs. The grey rows are real grey levels
  recovered from power alone; the binary rows are the best that design can produce,
  because it only ever sees a thresholded pixel.}
  \label{fig:grey}
\end{figure}

\begin{table}[h]
  \centering
  \small
  \resizebox{\textwidth}{!}{\begin{tabular}{lrrrrrr}
\hline
Design & MAE (0--255) & Corr. & MSSIM & F1 & Recog. & $n$ \\
\hline
Grey design (real 0--255 pixels) & 9.55 & 0.945 & 0.842 & 0.890 & 0.98 & 300 \\
Binary design (thresholded pixels) & 16.76 & -- & 0.570 & 0.815 & 0.91 & 300 \\
\hline
\end{tabular}
}
  \caption{The same generator, the same split, the same number of images; only the
  hardware design changes. MAE and correlation are between the recovered and the true
  grey levels.}
  \label{tab:grey}
\end{table}

On these numbers the grey design is better on every measure, which I did not expect.
The mean absolute error is $9.55$ on the 0--255 scale, under 4\,\% of full scale, the
correlation with the true grey level is $0.945$, MSSIM rises from $0.570$ to $0.842$ and
the recogniser reads $98\,\%$ of the recovered images correctly against $91\,\%$.

\textbf{This comparison is not controlled, and Section 8.6 shows the conclusion I drew
from it is wrong.} The grey bitstream differs from the binary one in two ways at once:
the input datapath is 8-bit instead of 1-bit, and its leakage amplifier is eight times
wider ($72\times63$ flops against $9\times63$). I attributed the whole improvement to the
first difference without testing the second. A third bitstream, built later, separates
them, and most of what is shown here turns out to be the amplifier. The paragraph that
follows in this subsection was written before that control existed; I have left it in
place so the correction in Section 8.6 has something to correct.

One caveat on the comparison of the MAE column: for the binary design the truth is 0 or
255, so an error is either 0 or 255 and the average is not measuring the same thing as it
is for the grey design. The correlation, MSSIM and recognition columns are the fair
comparisons, and they point the same way.

\subsection*{7.4 Repeating the generative reconstruction}

The Power2Picture route in the repository was trained on a capture made before the
authenticity gate existed. That capture has no design check and no functional check
recorded, and its bitstream hash does not match the verified one, so I treat its result
as void rather than comparing against it.

Repeating it on a gate-passing capture works. Training converges well before 60 epochs
(the 60 and 120 epoch runs restore the same best-validation checkpoint and score
identically), giving a foreground F1 of about $0.82$ and $0.91$ recognition.

\begin{table}[h]
  \centering
  \small
  \resizebox{\textwidth}{!}{\begin{tabular}{lrrrrrr}
\hline
Run & Pixel acc. & All-background & F1 & MSSIM & Recog. & $n$ \\
\hline
Earlier 9000-image run \emph{(void: no gate)} & 0.972 & 0.871 & 0.891 & 0.803 & 0.95 & 500 \\
Corrected capture, 20 epochs & 0.954 & 0.876 & 0.807 & 0.561 & 0.90 & 300 \\
Corrected capture, 60 epochs & 0.956 & 0.878 & 0.815 & 0.570 & 0.91 & 300 \\
Corrected capture, 120 epochs & 0.956 & 0.878 & 0.815 & 0.570 & 0.91 & 300 \\
\hline
\end{tabular}
}
  \caption{Generative reconstruction. The first row is the earlier result and is void:
  it was measured on a capture with no authenticity gate and a different bitstream.}
  \label{tab:p2p}
\end{table}

One caveat that should be stated with the number: this capture was made with a single
probe kernel and no averaging, which Section 6.1 shows is the weakest regime. So $0.82$
understates what the generative route can do, and it is not a fair comparison against
the template attack, which used five averaged repeats and nine kernels.


\section*{8. Follow-up work, 1 September}

Section 7 left four things open. Three of them are now measured and are reported here.
The fourth, the control bitstream for the quantisation question, is built and waits only
for a power cycle of the board.

I want to be honest about the shape of this section before the numbers start. Three
conclusions I wrote on 30 and 31 August did not survive these experiments. In each case
the reason was the same: I had drawn a rule from two measurements, and the rule broke as
soon as I measured a third point. I have left the wrong statements in the earlier
sections and marked them here, rather than editing them away, because the pattern is
itself the lesson.

\subsection*{8.1 Tuning the candidate search radius}

The candidate search keeps every template row whose power vector lies within a distance
$\delta$ of the measured one. I had been using $\delta = 1.0$ because that is the value
in the paper, and I never questioned it. Section 6.4 gave me a reason to: the 300-image
template scored \emph{worse} than the 40-image one, which is the wrong direction for more
data, and the candidate count per window had doubled. That is the signature of a radius
that is too wide for the template it is searching.

So I swept it. Both templates attack the same 40 evaluation images, only $\delta$ changes,
and the attack stage is the only thing that reruns because the captures already exist.
Thirteen values per template, twenty-seven runs, no board time.

\begin{table}[h]
  \centering
  \small
  \resizebox{\textwidth}{!}{\begin{tabular}{lrrrrrrrrr}
\hline
Template & $\delta$ = 1.4 & $\delta$ = 1 & $\delta$ = 0.6 & $\delta$ = 0.4 & $\delta$ = 0.25 & $\delta$ = 0.15 & $\delta$ = 0.08 & $\delta$ = 0.06 & $\delta$ = 0.03 \\
\hline
40 profiling images & 0.730 & 0.843 & 0.912 & 0.940 & \textbf{0.958} & 0.936 & 0.865 & 0.820 & 0.736 \\
300 profiling images & 0.696 & 0.759 & 0.849 & 0.890 & 0.926 & 0.949 & 0.967 & \textbf{0.969} & 0.946 \\
\hline
\end{tabular}
}
  \caption{Foreground F1 against the search radius. Both curves are unimodal and both
  peak far below the published value. The best value of each row is in bold.}
  \label{tab:delta}
\end{table}

\begin{figure}[h]
  \centering
  \includegraphics[width=\linewidth]{figures/fig15_delta_sweep.png}
  \caption{The same sweep drawn out. The dotted line is the value I had been using.}
  \label{fig:delta}
\end{figure}

The published $\delta$ was costing this project $0.13$ of foreground F1. With the radius
tuned, the 40-image template reaches $0.958$ and the 300-image template reaches
$\mathbf{0.969}$, with an MSSIM of $0.961$ and every one of the 40 digits recognised
correctly. The headline number of the whole project moves from $0.843$ to $0.969$.

Two things follow. The first is that more profiling data does help after all, but only if
the radius shrinks with it: $7.5$ times the template rows wanted roughly $4$ times the
tighter radius. At any fixed $\delta$ the larger template looks strictly worse, which is
exactly the trap I fell into in Section 6.4, and which anybody comparing reconstruction
quality against profiling-set size will fall into unless they tune the two together.

The second is a correction. In Section 6.4 I predicted the 300-image template would win
once $\delta$ was reduced. Then, with data only down to $\delta = 0.2$, I said it did not.
It does: the crossover is near $\delta = 0.1$, below where I had stopped looking. The
prediction was right and the retraction was wrong, and I only found out by sweeping
further than I thought was necessary.

\subsection*{8.2 Ruling out temperature}

Every clock-transfer number up to this point came from captures taken one condition after
the other: forty images at 5\,MHz, then forty at 10\,MHz. Glamo\v{c}anin et al. (DATE 2023)
show that this is the wrong way to acquire traces. Over a long acquisition the die warms
up, and if each class sits in its own block of time, a model can learn the thermal drift
between the blocks instead of the leakage. They also note the effect is strongest when the
leakage is weak, which describes my small convolution unit exactly. Their remedy is to
interleave the recordings so the drift is spread evenly.

This is a serious objection to Section 7.1, so I wrote a capture script that interleaves
the conditions per image: the same image is captured at both clocks back to back, then the
next image, and so on. The slow thermal ramp of the session then lands on both conditions
in the same proportion. Both conditions passed the functional check with $676/676$ outputs
bit-exact, so nothing else changed between the two acquisitions.

Comparing the two protocols on the same evaluation sets at $\delta = 0.3$, the largest
difference anywhere in the matrix is $0.039$ F1. The asymmetry that difference would have
to explain is $0.799$, so temperature accounts for at most about five per cent of the
effect. The direction is also reassuring: interleaving made the cross-clock case slightly
\emph{better}, meaning the sequential captures had been overstating the collapse rather
than inventing it.

I consider this the check that makes Section 7.1 reportable. Without it the whole result
rests on an acquisition protocol that a reviewer of that paper would reject immediately.

\subsection*{8.3 The transfer surface across four clocks}

Section 7.1 used two frequencies, and from those two points I wrote a rule: profile at the
higher clock, because it transfers downward for free. Two points are not enough to support
a rule like that, so I captured a proper grid.

The design is a four by four matrix. Forty profiling images (5300--5339) and forty
evaluation images (5200--5239), disjoint sets, each captured at 5, 7.5, 10 and 20\,MHz
with the conditions interleaved as in Section 8.2. That gives four templates and four
evaluation sets, and sixteen combinations. Every capture passed the functional check at
every clock, including 20\,MHz, which is four times the frequency the design was
originally built for.

\begin{table}[h]
  \centering
  \small
  \begin{tabular}{lrrrr}
\hline
Profiled at & \multicolumn{4}{c}{Attacked at} \\
 & 5\,MHz & 7.5\,MHz & 10\,MHz & 20\,MHz \\
\hline
5\,MHz & \textit{0.884} & 0.669 & 0.035 & 0.001 \\
7.5\,MHz & 0.901 & \textit{0.889} & 0.353 & 0.008 \\
10\,MHz & 0.848 & 0.893 & \textit{0.675} & 0.064 \\
20\,MHz & 0.026 & 0.162 & 0.305 & \textit{0.268} \\
\hline
\end{tabular}

  \caption{Foreground F1 for every combination of profiling and attack clock, at a fixed
  $\delta = 0.25$. The matched-condition diagonal is in italics.}
  \label{tab:grid}
\end{table}

\textbf{The rule I wrote from two points is wrong.} A template profiled at 20\,MHz is
almost useless at 5\,MHz, scoring $0.026$. Profiling at the highest available clock does
not transfer downward for free; it only looked that way because 5 and 10\,MHz happen to be
close enough to each other.

What the grid actually shows is a tolerance band, and the useful way to read it is by the
ratio between the two clocks rather than by their absolute values. The asymmetry I claimed
in Section 7.1 is real, and it holds in every pair I can form: for all six ordered pairs of
distinct clocks, profiling above and attacking below beats profiling below and attacking
above, and the margin grows with the ratio. At a ratio of two, $10 \to 5$ gives $0.851$
against $0.050$ for $5 \to 10$. At a ratio of four, $20 \to 5$ gives $0.071$ against
$0.019$ for $5 \to 20$, so by then both directions have failed and only the ordering
survives. Six pairs out of six is a stronger statement than the single ratio I quoted
before, and I have dropped the ``twenty times'' figure because it described one pair and
not the surface.

\subsubsection*{Tuning the radius for every cell}

The table above uses one $\delta = 0.25$ everywhere. That is fair between cells, but
Section 8.1 had already shown the best radius moves with the template, so a single value
cannot be each cell's optimum, and the cross-condition cells are a different problem again.
I therefore swept every cell over nine radii from $0.02$ to $0.60$, which is 144 attack
runs. No new capture was needed.

\begin{table}[h]
  \centering
  \small
  \begin{tabular}{lrrrr}
\hline
Profiled at & \multicolumn{4}{c}{Attacked at} \\
 & 5\,MHz & 7.5\,MHz & 10\,MHz & 20\,MHz \\
\hline
5\,MHz & \textit{0.884 {\tiny (0.25)}} & 0.706 {\tiny (0.15)} & 0.050 {\tiny (0.60)} & 0.019 {\tiny (0.60)} \\
7.5\,MHz & 0.901 {\tiny (0.25)} & \textit{0.918 {\tiny (0.15)}} & 0.510 {\tiny (0.10)} & 0.038 {\tiny (0.60)} \\
10\,MHz & 0.851 {\tiny (0.10)} & 0.895 {\tiny (0.15)} & \textit{0.882 {\tiny (0.06)}} & 0.132 {\tiny (0.60)} \\
20\,MHz & 0.071 {\tiny (0.60)} & 0.228 {\tiny (0.60)} & 0.305 {\tiny (0.25)} & \textit{0.345 {\tiny (0.60)}} \\
\hline
\end{tabular}

  \caption{The same surface with $\delta$ tuned for each cell, the chosen value in
  parentheses. The matched-condition diagonal is in italics. No cell peaks at either end
  of the swept range, so every optimum is bracketed.}
  \label{tab:gridbest}
\end{table}

The correction is large where I said it would be and small elsewhere. The 10\,MHz diagonal
moves from $0.675$ to $0.882$, because its best radius is $0.06$, four times tighter than
the value the fixed grid used. The other three diagonals move by less than $0.08$. The
off-diagonal cells barely move at all. So the shape of the surface was never in doubt, but
the fixed-$\delta$ table understated the matched-condition case badly enough that it should
not be quoted on its own.

One methodological point falls out of this. The radius is compared against a distance
between per-trace normalised feature vectors, and that distance does not have the same
scale at every clock: the median distance from an evaluation window to its nearest template
row is $0.18$ to $0.40$ when the evaluation clock is 20\,MHz and $0.65$ to $0.95$ when it is
5\,MHz, depending on which template is used. A single $\delta$ therefore
means a different acceptance rate in every cell, which is the real reason the fixed grid
was unfair to the diagonal. Tuning per cell is not a cosmetic improvement; it is what makes
the cells comparable at all.

\subsubsection*{The 20\,MHz row}

Earlier I wrote that the 20\,MHz row was weak and that I had no explanation. I can now say
more, though not everything.

First, how weak. The evaluation images are $11.35\,\%$ foreground, so a predictor that
marks every pixel as foreground reaches an F1 of $2p/(1+p) = 0.20$ while knowing nothing.
The 20\,MHz diagonal reaches $0.345$, which clears that floor, and its pixel-level
precision is $0.27$ against a base rate of $0.11$, so the reconstruction is more than twice
as likely to be right as chance wherever it marks foreground. The signal is real. But its
overall pixel accuracy is $0.794$ against $0.887$ for a blank image, because it
over-predicts foreground badly: recall $0.476$ against precision $0.270$. So the honest
description is a weak, badly calibrated reconstruction, not an absent one, and I should not
have called the row empty.

Second, why. The per-window leakage is measurably weaker at higher clocks. Correlating the
extracted feature against the true convolution output over twelve images gives $0.391$ at
5\,MHz, $0.224$ at 7.5, $0.123$ at 10 and $0.075$ at 20\,MHz, while the spread between the
five hardware repeats of the same window stays flat at about $0.34$. So this is not
measurement noise growing; the signal itself shrinks. Trace amplitude actually rises with
clock ($0.038$ to $0.051$ mean absolute), so it is not attenuation either.

I want to be careful about how much that table explains, because it is tempting to read
more into it than it supports. It does not explain the ordering of the other three clocks:
10\,MHz has a third of the correlation of 5\,MHz and an essentially equal diagonal
($0.882$ against $0.884$). The attack normalises each trace and matches on relative
distances, so an absolute correlation with the golden output is not the quantity it
consumes. What the table does support is the single narrow claim that 20\,MHz sits about
five times below the other three, and it is the only clock whose reconstruction falls below
the blank-image baseline.

Third, the confound, which is the reason I am not turning this into a countermeasure claim.
The ChipWhisperer-Lite samples at four times the FPGA clock, so the 20\,MHz capture runs the
ADC at 80\,MSa/s while the 5\,MHz capture runs it at 20. The sample count per window is
identical in all four captures --- $32$ samples, eight cycles, four samples per cycle, so
nothing is being averaged away --- but the sampling period is four times shorter, and the
analog front end has a fixed bandwidth. Weaker leakage at 20\,MHz and a less capable
instrument at 80\,MSa/s predict exactly the same measurement. I cannot separate them from
the captures I have. The clean discriminator is a 20\,MHz recapture at two samples per
cycle, which puts the ADC back at 40\,MSa/s with the accelerator unchanged; if the
correlation recovers, the instrument was the limit.

That matters because the obvious reading --- clock the accelerator faster and the side
channel goes away, which would be a far cheaper countermeasure than the clock randomisation
of the next subsection --- is exactly the reading the confound forbids. The accelerator
itself is fine at 20\,MHz: the functional check passes with $0$ mismatches on all $676$
outputs, so it computes correctly and only the measurement degrades. I am leaving this as an
open question with the experiment named, not as a result.

\subsection*{8.4 Does randomising the clock defend the accelerator?}

If a template is bound to the clock it was profiled at, then a defender who changes the
clock for every inference should break it. I suggested this as a cheap countermeasure,
cheaper than the dual-rail hiding or run-time bitstream swapping in the literature, so I
wanted to test it rather than assert it.

The grid of Section 8.3 makes the test almost free. Because every evaluation image was
captured at every clock, I can assemble the traces a clock-randomising defender would
produce instead of recapturing them: for each image, keep the trace from one randomly
chosen clock. The resulting set mixed 10, 7, 13 and 10 images across the four frequencies.
Its manifest records that it is synthetic and which clock each image came from, so it
cannot later be mistaken for a single-condition capture.

\begin{table}[h]
  \centering
  \small
  \begin{tabular}{lrrrrr}
\hline
Profiled at & Matched clock & Randomised clock & Change & Recognition & $\delta$ (m / r) \\
\hline
5\,MHz & 0.884 & 0.440 & -0.444 & 0.40 & 0.25 / 0.15 \\
7.5\,MHz & 0.918 & 0.601 & -0.317 & 0.57 & 0.15 / 0.10 \\
10\,MHz & 0.882 & 0.737 & -0.145 & 0.75 & 0.06 / 0.10 \\
20\,MHz & 0.345 & 0.248 & -0.097 & 0.30 & 0.60 / 0.25 \\
\hline
\end{tabular}

  \caption{What clock randomisation costs each template, compared with that template
  attacking its own matched clock.}
  \label{tab:jitter}
\end{table}

Both columns of this table are now tuned. When the matched-clock column was retuned per
cell in Section 8.3, leaving the randomised column at the old fixed $\delta = 0.25$ would
have measured the defence against a handicapped attacker, so the randomised set was swept
over the same nine radii and each row reports each side at its own best value.

The defence is much weaker than I expected. Against a template profiled at 5\,MHz it works
well, cutting F1 from $0.884$ to $0.440$. But against a template profiled at 10\,MHz it
costs $0.145$, from $0.882$ to $0.737$, and the recogniser still reads 75 per cent of the
digits.

Retuning did change this number and I should say so plainly: the earlier version of this
table reported a loss of only $0.018$ for the 10\,MHz template, and I described that as
achieving almost nothing. The true cost is eight times larger. But it does not rescue the
defence, because the absolute number moved the wrong way for the defender --- a
correctly-tuned attacker who profiled at 10\,MHz comes out of clock randomisation at
$0.737$ rather than the $0.656$ I reported, and reads three digits in four instead of two
in three.

The reason is visible in the grid. A template profiled at 10\,MHz sits in the middle of
the frequency range, within transfer distance of 5, 7.5 and 10\,MHz alike, so it handles
most of the randomised images no matter which clock they came from. An attacker does not
need to know the defender's schedule; he only needs to profile somewhere central.

The conclusion is more specific than ``clock randomisation does not work''. It works only
if the defender randomises \emph{outside the attacker's tolerance band}, and Section 8.3
measured that band as roughly a factor of two. Varying the clock between 5 and 10\,MHz is
not enough. A defender would have to span a much wider range, or hop between frequencies
far enough apart that no single profiling point covers them, and that is a considerably
more expensive thing to build than I implied when I proposed it.

\subsection*{8.5 Mixing conditions, with the coverage error fixed}

Section 7.1 tested a mixed template built from 20 images at 5\,MHz and 20 at 10\,MHz, and
I concluded that mixing helps but loses to profiling at the higher clock. That experiment
had a flaw I found afterwards: the mixed template held only 20 images of coverage per
condition, against 40 for each pure template. So it varied two things at once, the mixing
and the per-condition coverage, and I could not attribute the difference to either.

The repeat holds coverage constant. Each pure template uses 40 images at its own clock;
the mixed template uses 40 at each clock, so it has the same per-condition coverage as the
pure templates and twice as many rows. Because Section 8.1 showed the best radius depends
on template size, each template is swept over four values of $\delta$ rather than compared
at one.

\begin{table}[h]
  \centering
  \small
  \begin{tabular}{llrr}
\hline
Template & $\delta$ & Attacked at 5\,MHz & Attacked at 10\,MHz \\
\hline
Profiled at 5\,MHz & 0.02 & 0.534 & 0.036 \\
 & 0.03 & 0.581 & 0.036 \\
 & 0.04 & 0.640 & 0.029 \\
 & 0.06 & 0.728 & 0.035 \\
 & 0.10 & 0.835 & 0.069 \\
 & 0.15 & 0.878 & 0.065 \\
 & 0.25 & 0.891 & 0.077 \\
 & 0.40 & 0.850 & 0.061 \\
\hline
Profiled at 10\,MHz & 0.02 & 0.821 & 0.674 \\
 & 0.03 & 0.831 & 0.790 \\
 & 0.04 & 0.835 & 0.847 \\
 & 0.06 & 0.828 & 0.875 \\
 & 0.10 & 0.837 & 0.872 \\
 & 0.15 & 0.837 & 0.813 \\
 & 0.25 & 0.846 & 0.687 \\
 & 0.40 & 0.833 & 0.525 \\
\hline
Mixed, 40 at each & 0.02 & 0.627 & 0.305 \\
 & 0.03 & 0.679 & 0.353 \\
 & 0.04 & 0.734 & 0.404 \\
 & 0.06 & 0.810 & 0.429 \\
 & 0.10 & 0.871 & 0.469 \\
 & 0.15 & 0.896 & 0.369 \\
 & 0.25 & 0.885 & 0.306 \\
 & 0.40 & 0.835 & 0.206 \\
\hline
\end{tabular}

  \caption{Mixing with equal per-condition coverage. The mixed template has twice the rows
  of the pure ones, so each is swept over $\delta$ rather than compared at a single value.}
  \label{tab:mix2}
\end{table}

Mixing still loses at 10\,MHz, and not narrowly: the mixed template reaches $0.469$ while
the pure 10\,MHz template reaches $0.875$. Since the two now hold the same number of
10\,MHz images, the 5\,MHz rows are not diluting coverage; they are actively supplying
wrong candidates that survive the distance test.

That last sentence is the mechanism I proposed on 31 August and then withdrew. Withdrawing
it was a mistake caused by the coverage flaw. In the flawed experiment, tightening $\delta$
made mixing worse, which contradicted the mechanism; with coverage fixed, tightening
$\delta$ from $0.40$ to $0.10$ makes mixing better ($0.206$ rising to $0.469$), which is
what the mechanism predicts.

I then made the same bracketing error here that Section 8.3 corrects in the grid. This
sweep originally stopped at $\delta = 0.10$, and the mixed template's best value sat on
that floor, so I wrote that tightening $\delta$ helps mixing and left it there. Taken
literally my own mechanism predicts the improvement should continue below $0.10$, and that
prediction is wrong: extending the sweep to $0.02$ turns the curve over, from $0.469$ at
$0.10$ down to $0.305$ at $0.02$. So the mechanism is only half the story. Tightening the
radius does remove wrong-condition candidates, which helps, but past a point it also starts
removing the right ones, which hurts, and the mixed template's optimum is where those two
effects balance. That is ordinary radius tuning rather than anything special about mixing,
and the honest version of the claim is just that the mixed template's ceiling ($0.469$) is
far below the pure one's ($0.875$) at every radius either of them prefers.

Extending the sweep also turned up something I had not seen. At 5\,MHz mixing is not worse
at all: the mixed template reaches $0.896$ at $\delta = 0.15$ against $0.891$ for the pure
5\,MHz template. The two are within a few thousandths of each other, which on 40 images I
would call a tie rather than a win, but it means ``mixing does not help'' is only true at
10\,MHz. Where mixing hurts, it hurts the harder condition.

The practical answer turns out to be simpler than mixing. A template profiled at 10\,MHz
scores $0.837$ at 5\,MHz with $\delta = 0.10$ and $0.875$ at 10\,MHz with $\delta = 0.06$. That beats the mixed
template at both clocks while using half the profiling data. Given Section 8.3, I would now
state the guidance as: profile at a frequency near the middle of the range you expect to
attack, and tune $\delta$ for the size of the template you built. Do not mix.


\subsection*{8.6 Was it the input width, or just a bigger amplifier?}

Section 7.3 compared the grey design against the binary one and concluded that removing
the input binarisation made the attack better. That comparison had a flaw I did not see
at the time. The two bitstreams differ in two ways, not one. The grey design takes an
8-bit pixel instead of a 1-bit pixel, and its leakage amplifier is eight times wider,
$72\times63 = 4536$ flops against $9\times63 = 567$, because the amplifier rides a
byte-wide term instead of a bit-wide one. A wider bank toggling on every window produces
a larger power signal regardless of what the input carries, so the experiment could not
distinguish ``the input carries more information'' from ``the circuit makes more noise''.

The control is a third bitstream. It keeps the binary datapath unchanged --- one bit per
pixel, the same 9-bit match term, the same convolution --- and only widens the amplifier
to 72 bits by replicating the same nine bits eight times. Replication adds amplitude but
no information, so if the amplifier is what mattered, this design should behave like the
grey one; if the input is what mattered, it should behave like the original binary one.
It passes the same simulation as the binary core, bit-exact on all 676 outputs.

That the wide design was really running is worth stating, because a readback cannot tell
the two binary bitstreams apart --- they have the same register interface. The evidence is
physical: on the same images and the same settings, the mean absolute trace value rises
from $1341$ to $6339$, a factor of $4.7$. Nothing but the widened bank can do that.

I also found and fixed a second flaw while setting this up. The grey and binary captures
of Section 7.3 used different images (7000--8999 against 2000--3999), so they were not
comparable even before the amplifier question. All three designs were therefore captured
again over the same 2000 images, with the same shard size, kernel, threshold and shuffle
seed.

\begin{table}[h]
  \centering
  \small
  \resizebox{\textwidth}{!}{\begin{tabular}{lrrrrr}
\hline
Design & MAE (0--255) & Corr. & MSSIM & F1 & Recog. \\
\hline
Binary, narrow amplifier ($9\times63$) & 16.76 & -- & 0.570 & 0.815 & 0.91 \\
Binary, wide amplifier ($72\times63$) & 11.85 & 0.892 & 0.662 & 0.878 & 0.97 \\
Grey, wide amplifier ($72\times63$) & 9.76 & 0.929 & 0.833 & 0.855 & 0.95 \\
\hline
\end{tabular}
}
  \caption{The amplifier control. All three rows use the same 2000 images and the same
  generator; the first two rows also use the same task, and differ only in amplifier
  width.}
  \label{tab:amp}
\end{table}

The result changes the conclusion. Widening the amplifier alone, with the binary datapath
untouched, moves F1 from $0.815$ to $0.878$ and recognition from $0.91$ to $0.97$. The
grey design then scores $0.855$ and $0.95$, which is \emph{worse} than the wide binary
design on both. So the sentence in Section 7.3, that the grey design is better on every
comparable measure, is false: it was true only against the narrow amplifier, and the
amplifier is not what I was trying to study.

What does survive is a narrower and, I think, more interesting claim. Grey wins clearly on
the measures that describe intensity rather than shape: MSSIM $0.833$ against $0.662$,
which is the largest single change in the table, and correlation $0.929$ against $0.892$.
It achieves that while solving a harder problem, because the grey target is a continuous
0--255 image and the binary target is a 1-bit silhouette. F1 counts foreground overlap and
mildly penalises the harder task, which is why the ordering flips between the two groups
of columns.

The honest statement is therefore: removing the input binarisation does not improve
silhouette recovery, a wider leakage amplifier does that. What it buys is the ability to
recover pixel intensity at all, which the binary datapath cannot do by construction. That
is a smaller claim than the one in Section 7.3, and it is the one I can defend.

One limitation remains. The wide-against-narrow comparison is fully controlled: same task,
same images, one variable. The grey-against-wide comparison now has matched images and a
matched amplifier, but the two designs are still solving different reconstruction problems,
so I would not put a single number on how much the input width is worth.


\subsection*{8.7 Where cross-condition transfer actually breaks}

Having tuned the radius everywhere, I wanted to know why the upward cells fail, because
``the template is bound to its clock'' describes the result without explaining it. The
attack has three stages that can each lose the answer, and they can be measured separately.
For each convolution window the nine probe features are split into three groups of three;
each group retrieves every template row within radius $\sqrt{\delta}$ of the query in its
own three-dimensional subspace; the three sets of candidate 3$\times$3 patches are
intersected; and the surviving candidates are handed to the greedy overlap search that
assembles the image. I instrumented the first two stages using the real template files and
the real evaluation bundles, so the numbers below describe exactly what the attack sees.

\begin{table}[h]
  \centering
  \small
  \begin{tabular}{lrrrr}
\hline
Profiled $\to$ attacked & In set & Rank 1 & Candidates & F1 \\
\hline
\textit{5 $\to$ 5\,MHz} & 0.836 & 0.760 & 4.0 & 0.884 \\
5 $\to$ 7.5\,MHz & 0.795 & 0.699 & 4.7 & 0.669 \\
5 $\to$ 10\,MHz & 0.778 & 0.681 & 5.3 & 0.035 \\
5 $\to$ 20\,MHz & 0.737 & 0.688 & 6.1 & 0.001 \\
\hline
7.5 $\to$ 5\,MHz & 0.735 & 0.672 & 3.8 & 0.901 \\
\textit{7.5 $\to$ 7.5\,MHz} & 0.867 & 0.759 & 5.6 & 0.889 \\
7.5 $\to$ 10\,MHz & 0.865 & 0.724 & 7.8 & 0.353 \\
7.5 $\to$ 20\,MHz & 0.768 & 0.630 & 8.7 & 0.008 \\
\hline
10 $\to$ 5\,MHz & 0.593 & 0.500 & 5.9 & 0.848 \\
10 $\to$ 7.5\,MHz & 0.793 & 0.677 & 8.4 & 0.893 \\
\textit{10 $\to$ 10\,MHz} & 0.925 & 0.744 & 15.5 & 0.675 \\
10 $\to$ 20\,MHz & 0.853 & 0.622 & 21.4 & 0.064 \\
\hline
20 $\to$ 5\,MHz & 0.499 & 0.328 & 14.8 & 0.026 \\
20 $\to$ 7.5\,MHz & 0.606 & 0.406 & 23.9 & 0.162 \\
20 $\to$ 10\,MHz & 0.799 & 0.533 & 35.4 & 0.305 \\
\textit{20 $\to$ 20\,MHz} & 0.955 & 0.655 & 56.3 & 0.268 \\
\hline
\end{tabular}

  \caption{Candidate quality against final result, at $\delta = 0.25$. ``In set'' is how
  often the true patch survives the three-way intersection; ``rank 1'' is how often it is
  also the closest candidate in the full nine-dimensional feature space.}
  \label{tab:stage}
\end{table}

The result is not what I expected, and it points somewhere I had not been looking. Take the
two cells of the 5 and 10\,MHz pair. Profiling at 5 and attacking at 10 puts the true patch
in the candidate set $77.8\,\%$ of the time and ranks it first $68.1\,\%$ of the time --- and
the finished reconstruction scores $0.035$. Profiling at 10 and attacking at 5 does
\emph{worse} at both: the true patch is in the set $59.3\,\%$ of the time and ranks first
$50.0\,\%$ --- and the finished reconstruction scores $0.848$. The cell with the better
candidates produces the far worse image, and the ordering is inverted for the whole pair.

My first worry was that this was an artefact of the radius, since the table above is taken
at $\delta = 0.25$ and neither of those two cells actually prefers that value. So I repeated
the diagnostic at $0.10$ and at $0.60$, which bracket both cells' own optima. The inversion
holds at all three radii: at $0.10$ the pair reads $70.7\,\%$ in set and $66.9\,\%$ rank one
upward against $42.7\,\%$ and $38.3\,\%$ downward, and at $0.60$ it reads $84.8\,\%$ and
$67.9\,\%$ upward against $72.5\,\%$ and $55.8\,\%$ downward. The sharpest way to put it is
to compare each cell at the radius it actually wins on: profiling at 5 and attacking at 10
retrieves the true patch $84.8\,\%$ of the time and ranks it first $67.9\,\%$ of the time,
and scores $0.050$; profiling at 10 and attacking at 5 retrieves it only $42.7\,\%$ of the
time and ranks it first $38.3\,\%$, and scores $0.851$. Twice the retrieval quality, a
seventeenfold worse image.

So for this pair the asymmetry does not live in the template match. Retrieval and ranking
are better in the direction that fails, at every radius I tried. Whatever destroys these
upward cells happens in the greedy assembly stage, which is the one part of the pipeline
that does not look at each window independently: it places patches by consistency with the
pixels already placed, so it needs the candidate sets of neighbouring windows to agree with
each other, not merely to be individually correct. A systematic offset between conditions
can leave every window's candidates good while making neighbouring windows disagree, and
that is the shape of failure the numbers point at.

I have not proved that mechanism, and I am not going to assert it on the strength of one
diagnostic. I also want to be careful about how far the negative result reaches, because
the same table shows retrieval still matters elsewhere in the grid: along the 20\,MHz row,
where the leakage itself is weak, the true patch is in the candidate set $49.9\,\%$ of the
time against 5\,MHz and $79.9\,\%$ against 10\,MHz, and the reconstruction follows in the
same order. So retrieval is clearly still binding there. The claim I am willing to defend
is the narrow one: in the 5 and 10\,MHz pair, cross-condition failure is not caused by the
power template failing to match, so better features, better normalisation or a better
distance would not fix those two cells --- those stages are already working better in the
cell that fails than in the cell that succeeds. That is a more specific research question
than the one I started with, and it is where I would look next.


\section*{9. Follow-up work, 6 September: a second dataset}

Everything up to here uses MNIST. Wei et al. state that their recovery template is
``independent from the chosen dataset if the template is built in a way avoids
overfitting'', on the grounds that it maps power to the pixels inside a convolution
window rather than to anything semantic. The evidence they give is one mammography
image, shown as a picture with no numbers, recovered with the MNIST template; they list
harder datasets as future work. Power2Picture did evaluate Fashion-MNIST, but only for
the generative attack, and it degraded sharply there: MSSIM $0.62 \to 0.38$, pixel
distance $32 \to 77$, recognition $96.4\,\% \to 53.0\,\%$. Nobody has run a template
attack on Fashion-MNIST. So the claim that matters most for reusing a template is the
one with the least evidence behind it, and I can test it directly.

Fashion-MNIST costs almost nothing to add. The images are $28\times28$ 8-bit, exactly the
geometry the FPGA already takes, so no design change was needed and the capture script
only needed a different image file. What did change is the statistics, and they change
enough to break the metrics if it is not handled.

\subsection*{9.1 Why the usual numbers stop being comparable}

MNIST is sparse and Fashion is not. On the evaluation sets used here MNIST is
$11.4\,\%$ foreground and Fashion is $25.9\,\%$. That moves two baselines at once. A
blank image scores $0.886$ pixel accuracy on MNIST and only $0.741$ on Fashion, and the
best F1 reachable by a predictor that knows nothing at all --- mark every pixel
foreground, giving precision $p$ and recall $1$ --- rises from $2p/(1+p) = 0.204$ to
$0.474$.

So a raw F1 of $0.55$ would be a decent result on MNIST and almost nothing on Fashion,
and a \emph{lower} pixel accuracy on Fashion can still be a better reconstruction. Every
metric in this section is therefore reported three ways: the raw value, the chance
ceiling for that particular truth set, and the excess, $(\mathrm{F1} -
\mathrm{chance})/(1 - \mathrm{chance})$, which is $0$ for a label-blind predictor and $1$
for a perfect one. The chance figure is computed from the truth images actually passed to
the scorer, not from the dataset as a whole, because the 40-image evaluation subset
($25.9\,\%$) is not the same as the full Fashion test set ($31.5\,\%$).

I would not have caught the size of this effect without the excess column, and the first
version of this experiment was scored with the wrong recogniser as well --- the scoring
path builds the digit classifier with no argument, so recovered clothing was being read
by an MNIST model, which reported $0.025$ accuracy on the ground-truth Fashion images
where the correct classifier gets $0.842$. A Fashion recogniser was trained for this
section and reaches $0.902$ on greyscale and $0.842$ on binarised images, against
$0.172$ on MNIST digits, which is near the $0.10$ floor and confirms it is really a
clothing model.

\subsection*{9.2 The dataset transfer grid}

One cross run cannot answer the question, because a poor MNIST-to-Fashion result would
confound ``the template did not transfer'' with ``Fashion is simply harder''. The
control is the same grid structure used for the clocks in Section 8.3: two templates,
two evaluation sets, four cells, with $\delta$ swept over twelve radii per cell and every
optimum confirmed to be interior. Both datasets were captured at 10\,MHz with identical
gain, averaging, dwell, threshold, kernel and shuffle seed, so the dataset is the only
variable. The MNIST row reuses the existing 10\,MHz cell, and reproduces its diagonal
exactly at $0.882$, which is the consistency check rather than a new measurement.

\begin{table}[h]
  \centering
  \small
  \begin{tabular}{llrrrrr}
\hline
Template & Attacked & $\delta$ & F1 & chance & excess & MSSIM \\
\hline
MNIST & MNIST & 0.060 & 0.882 & 0.204 & 0.852 & 0.852 \\
MNIST & Fashion & 0.060 & 0.432 & 0.412 & 0.034 & 0.432 \\
Fashion & MNIST & 0.040 & 0.810 & 0.204 & 0.762 & 0.761 \\
Fashion & Fashion & 0.020 & 0.801 & 0.412 & 0.662 & 0.626 \\
\hline
\end{tabular}

  \caption{Template dataset against attacked dataset, $\delta$ tuned per cell. ``Excess''
  rescales F1 so that $0$ is a label-blind predictor and $1$ is perfect; it is the only
  column comparable across the two datasets.}
  \label{tab:dsgrid}
\end{table}

Two things fall out. Fashion-MNIST \emph{is} recoverable from power: the matched cell
reaches $0.662$ excess, and the recogniser reads $72\,\%$ of the recovered garments.
That is the first template-attack result on this dataset. But the MNIST template on
Fashion traces reaches $0.034$ excess, which is chance. Raw F1 for that cell is $0.432$
and looks like partial success; against a chance ceiling of $0.412$ it is nothing. So
the template does not transfer, and because the control on the same traces works, that
is a property of the template rather than of the difficulty of the images.

\subsection*{9.3 Why it fails, and the condition under which it would work}

The transfer is also asymmetric. Fashion to MNIST reaches $0.762$ excess, close to
MNIST's own $0.852$, while MNIST to Fashion collapses. That asymmetry points at the
answer, because the template is not really a power-to-pixel function at all. It is a
dictionary: a list of observed power vectors, each labelled with the $3\times3$ patch
that produced it. Its vocabulary is inherited from whatever images were profiled.

\begin{table}[h]
  \centering
  \small
  \begin{tabular}{lrrrr}
\hline
Template & distinct codes & all-zero & all-ones & mean density \\
\hline
MNIST & 156\,/\,512 & 0.683 & 0.032 & 0.169 \\
Fashion & 450\,/\,512 & 0.489 & 0.190 & 0.351 \\
\hline
\end{tabular}

  \caption{What each template actually contains. Both hold $27\,040$ rows from 40 images;
  they differ in how many of the 512 possible $3\times3$ binary patches ever occurred.}
  \label{tab:patchcov}
\end{table}

The MNIST template contains 156 of the 512 possible patches and the Fashion one 450. Of
these, 154 are shared, 296 occur only in the Fashion template, and just 2 only in the
MNIST one. The MNIST template is very nearly a \emph{subset} of the Fashion template. It
cannot emit a dense patch because it never observed one, so on clothing it under-predicts
foreground by a factor of two and a half --- $0.079$ against $0.196$ true --- and solid
garments dissolve.

\begin{figure}[h]
  \centering
  \includegraphics[width=\textwidth]{figures/fig_dataset_transfer.png}
  \caption{Fashion-MNIST recovered by the template attack. Top: truth. Middle: the
  Fashion-built template. Bottom: the MNIST-built template, which renders the thin
  shoes but erases the solid garments.}
  \label{fig:dstransfer}
\end{figure}

Per class the effect is monotone in density. With the MNIST template, sneakers
($14.5\,\%$ foreground) come back with $93\,\%$ of their foreground intact and classify
perfectly, while a coat ($53.3\,\%$ foreground) comes back with $2\,\%$ of its foreground
and never classifies correctly. Across the ten classes the correlation between true
foreground density and the fraction recovered is $-0.77$ for the MNIST template against
$-0.52$ for the Fashion control. Recognition splits the same way: $0.762$ on the thin
classes and $0.158$ on the solid ones. That also resolves an apparent contradiction, since
the overall recognition of $0.475$ looks incompatible with an F1 excess of $0.034$ until
one sees that the successes are concentrated in the sparse classes that F1 weights least.

The useful conclusion is not that the paper is wrong. It is more specific, and it
explains their result rather than contradicting it: the template is a patch dictionary,
and transfer works when the target's patch set is covered by the profiling set. That is
a directional condition, which is why Fashion transfers to MNIST and not the reverse. It
also fits their mammography demonstration, since a radiograph is largely smooth gradient
and therefore low-coverage --- the direction that works. Anyone reusing a profiled
template should check patch coverage first, and it costs nothing to check.

\subsection*{9.4 Greyscale clothing}

Everything above is binary, because the template attack classifies 1-bit $3\times3$
patches and cannot express a grey level by construction. That is a poor fit for this
dataset: $25.8\,\%$ of Fashion pixels are mid-tone against $5.4\,\%$ for MNIST, so
binarising discards a quarter of the image instead of a twentieth. Recovering clothing
properly needs the grey design and the generative route.

\begin{table}[h]
  \centering
  \small
  \begin{tabular}{lrrrrrr}
\hline
Run & MAE (0--255) & Corr. & MSSIM & F1 & Recog. & (original) \\
\hline
MNIST, grey & 9.76 & 0.929 & 0.833 & 0.855 & 0.970 & 0.980 \\
Fashion, grey & 20.54 & 0.910 & 0.737 & 0.820 & 0.827 & 0.830 \\
\hline
\end{tabular}

  \caption{Greyscale reconstruction with identical generator settings, features and
  split. ``Recog.'' is on the greyscale reconstruction, with the accuracy on the
  original images beside it as the ceiling.}
  \label{tab:fashgrey}
\end{table}

\begin{figure}[h]
  \centering
  \includegraphics[width=\textwidth]{figures/fig_fashion_grey.png}
  \caption{Fashion-MNIST greyscale recovered from power alone.}
  \label{fig:fashgrey}
\end{figure}

The reconstruction carries real intensity rather than a silhouette: the recovered pixel
spread is $0.324$ against $0.350$ in the truth, and the shading on a pullover's sleeves
or a bag's clasp is visible in the figure. The strongest number is the recognition
column. A Fashion classifier reads the power-recovered images at $0.827$ against $0.830$
on the originals, so in classification terms this reconstruction is very nearly lossless,
even though the mean absolute error of $20.5$ on the 0--255 scale is twice MNIST's.
Part of that error gap is density rather than difficulty --- Fashion simply has more
non-background pixels to get wrong --- and the correlation, $0.910$ against $0.929$,
is the fairer comparison and shows a much smaller gap.

These numbers are better than the published Fashion-MNIST figures in Power2Picture, but
I do not present that as an improvement on their work and it should not be read as one.
They attack a full LeNet through on-chip TDC sensors on a PYNQ-Z1. I attack a single
convolution layer through an external shunt on a CW305, in a design that contains a
leakage amplifier I added specifically to make it leak. My target is far simpler and
deliberately instrumented in the attacker's favour, so the two sets of numbers are not
measuring the same difficulty. What this does establish is that the grey pipeline works
on a dense, non-digit dataset, which was untested.

\subsection*{9.5 How much of this is the accelerator, and how much is the amplifier?}

Every number above comes from a design carrying a leakage amplifier I added: a bank of
registers whose only purpose is to toggle on the MAC term so the core leaks harder. At
the width used throughout the report that bank is $72\times63 = 4536$ flops. Neither
Wei et al. nor Power2Picture add such a circuit, so a reader is entitled to ask whether
these results describe an attack on a neural accelerator or an attack on a test bench
built to be attacked. This subsection puts a number on that, and it is the number I would
want a reviewer to see first.

The amplifier width is a parameter, so the sweep is four builds of the same design with
nothing else changed. The synthesis reports confirm the parameter does what it claims:
register count over the unamplified build rises by exactly $72\times$lanes, so $72$,
$216$, $1080$ and $4536$ flops. Each was captured over the same 2000 Fashion images and
reconstructed with the same generator, features, split and seed.

\begin{table}[h]
  \centering
  \small
  \resizebox{\textwidth}{!}{\begin{tabular}{rrrrrrrr}
\hline
Lanes & Amp.\ flops & Mean $|$trace$|$ & MAE & Corr. & MSSIM & F1 excess & Recog. \\
\hline
63 & 4536 & 7602 & 20.97 & 0.909 & 0.730 & 0.650 & 0.857 \\
15 & 1080 & 3314 & 25.88 & 0.872 & 0.650 & 0.579 & 0.753 \\
3 & 216 & 2229 & 36.60 & 0.759 & 0.511 & 0.399 & 0.443 \\
1 & 72 & 2197 & 37.26 & 0.750 & 0.503 & 0.422 & 0.400 \\
\hline
\end{tabular}
}
  \caption{Greyscale reconstruction against leakage-amplifier size. The recogniser
  ceiling on the original images is $0.830$ and chance is $0.100$.}
  \label{tab:ampsweep}
\end{table}

\begin{figure}[h]
  \centering
  \includegraphics[width=\textwidth]{figures/fig_amplifier_sweep.png}
  \caption{The same test images recovered at each amplifier width.}
  \label{fig:ampsweep}
\end{figure}

The answer is not the one I expected. Cutting the amplifier by a factor of $63$, from
$4536$ flops to $72$, costs a great deal --- recognition falls from $0.857$ to $0.400$ and
MSSIM from $0.730$ to $0.503$ --- but the attack plainly still works. At $72$ flops the
correlation between the recovered and the true grey level is $0.750$ and the recogniser
reads four times as many garments as chance.

More telling is the shape. The last step cuts the bank by a further factor of three and
changes almost nothing: correlation $0.759 \to 0.750$, MSSIM $0.511 \to 0.503$, and F1
excess actually rises. The measured trace amplitude behaves the same way, $2229 \to 2197$,
a difference of $1.4\,\%$, because by this point the amplifier contributes far less power
than the rest of the design and the total is dominated by the register file, the USB
interface and the clocking. The curve has flattened well before the smallest bank in this
sweep, and Section 9.7 confirms it continues flat to zero. The smallest bank I can
build, which says the leakage the attack is living on at that end is coming from the
convolution datapath itself rather than from anything I added.

That corrects something I had written down as settled. My working note from August said
the attack has no signal without the amplifier. On the binary design that may still hold;
on the grey design it is wrong, and most of the amplifier's contribution is already spent
by $216$ flops.

One caveat governs everything in this subsection, and it is important enough that the
next one is devoted to it. Every point above was measured with the same small profiling
set of about $1500$ traces. That makes the comparison between widths fair, but it fixes
the other quantity an attacker controls, and Section 9.6 shows the degradation here is
not the hard limit it looks like: most of it can be bought back with more profiling
data. Read this table as the cost of a narrower amplifier \emph{at a fixed and rather
small profiling budget}, not as the cost in general.

The zero point is now measured, and Section 9.7 gives it. It cost three failed builds to
get there, and the reason is worth stating here because it invalidated a claim this report
previously made. Two earlier attempts at a bank-free build produced a design that
programmed cleanly, reported DONE high, and answered zero at every register address, and I
wrote that I had no explanation for it. The explanation was mine: with the bank gone I had
left the leakage output driven combinationally, \texttt{assign leak\_sink = \^{}conv\_value},
which runs a path from distributed RAM through nine $16$-bit adders straight to a top-level
pad. Every working build in this report drives that pin from a register. One flop restores
it, and one flop is not an amplifier --- the smallest bank in the sweep above is $72$.

\subsubsection*{The check that made this trustworthy}

The first attempt at this sweep produced a clean-looking curve that was entirely false,
and the way it was caught is worth recording. All four bitstreams are grey designs with
identical register maps, so neither the stock-AES page signature nor the byte-ramp probe
distinguishes them, and this board does not reprogram over an already-configured FPGA:
the call reports success, DONE stays high, and the previous design keeps running. The
four captures came back at $7471$, $7478$, $7486$ and $7500$ mean absolute amplitude, a
spread of $0.4\,\%$ in the wrong direction, with identical functional signatures. A design
with $4536$ toggling flops and one with none cannot produce the same trace. They were four
captures of one bitstream.

Had I not fixed an amplitude criterion in advance, that run would have gone into this
report as a smooth curve of reconstruction quality against amplifier width, and every
point in it would have been the same measurement. The correct procedure, which produced
the table above, is: power-cycle the board, program while it is blank, confirm the
register interface responds, confirm the functional check is bit-exact on all 676 outputs,
and only then capture. That is encoded in \texttt{host/program\_and\_capture.py}, which
refuses to capture if the FPGA is already configured or if the interface does not answer.

One provenance note. The $216$-flop point was built from a slightly different revision of
the core than the other three, because a rebuild batch was interrupted before it. It
passes both gates and sits where the curve predicts, so I have kept it, but it is the one
row not built from the same source as its neighbours.


\subsection*{9.6 Running it the way the paper does}

Section 9.5 leaves an obvious hole. It varies the amplifier while holding the profiling
set at about $1500$ traces, which is nothing like the $60\,000$ Power2Picture train on,
so a narrower amplifier and a smaller profiling set are confounded. The way to separate
them is to run the attack the way the paper actually runs it and then vary both.

Reproducing their training side turned out to need less work than I expected, because the
generator in this repository was already built to their Table I: three fully connected
layers of $128$, $128$ and $12544$, reshaped to $256\times7\times7$, then three transposed
convolutions of $128$, $64$ and $1$ features. Adam at $0.001$ matched too. What did not
match was the loss, the batch size and, most of all, the amount of data.

\begin{table}[h]
  \centering
  \small
  \resizebox{\textwidth}{!}{\begin{tabular}{lll}
\hline
Aspect & Power2Picture & This work \\
\hline
Generator architecture & Table I: 128, 128, 12544, 3 transposed convs & identical \\
Optimiser, learning rate & Adam, $0.001$ & identical \\
Loss & \texttt{gradmse} (MSE + gradient MSE) & identical \\
Batch size & 256 & identical \\
Profiling set & 60k train split & identical \\
Dataset & Fashion-MNIST & identical \\
Victim network & quantised LeNet, full & one $3\times3$ conv layer \\
Measurement & 5 on-chip TDC sensors & external shunt, CW-Lite \\
Window rate & streaming, one window per cycle & $2$ cycles per window (\texttt{DWELL}), min. supported \\
Leakage amplifier & none & none ($1$ flop output register) \\
\hline
\end{tabular}
}
  \caption{Where this setup now matches Power2Picture and where it cannot.}
  \label{tab:papersettings}
\end{table}

The two settings I could change cheaply, I changed first, and they went the wrong way. On
the $1500$-trace captures, switching from MSE to their \texttt{gradmse} loss cost MSSIM
($0.730 \to 0.705$) and moving the batch from $128$ to $256$ cost more ($\to 0.681$). The
paper reports that \texttt{gradmse} significantly improves recovery, so this looked like a
contradiction until the obvious explanation: both settings are chosen for a $60\,000$-image
training run, and a batch of $256$ over $1500$ samples is six gradient steps per epoch.
The settings are coupled to the data volume, and testing them at the wrong volume tests
nothing.

So I captured the real thing: the full $60\,000$-image Fashion training split, at both the
smallest and the largest amplifier. That is $1.6$\,GB of traces per width and about
twenty-five minutes of capture each, both passing the functional check bit-exact on all
$676$ outputs, and with mean trace amplitudes of $2187$ and $7646$ confirming the two
bitstreams really were different.

\begin{table}[h]
  \centering
  \small
  \resizebox{\textwidth}{!}{\begin{tabular}{lrrrrrr}
\hline
Amplifier, profiling set & MAE & Corr. & MSSIM & F1 excess & Recog. & (ceiling) \\
\hline
none ($1$ flop), $56$k traces & 23.98 & 0.878 & 0.724 & 0.632 & 0.746 & 0.898 \\
$72$ flops, $1.5$k traces & 37.26 & 0.750 & 0.503 & 0.422 & 0.400 & 0.830 \\
$72$ flops, $56$k traces & 23.75 & 0.880 & 0.725 & 0.635 & 0.757 & 0.898 \\
$4536$ flops, $1.5$k traces & 20.97 & 0.909 & 0.730 & 0.650 & 0.857 & 0.830 \\
$4536$ flops, $56$k traces & 14.49 & 0.950 & 0.861 & 0.759 & 0.919 & 0.898 \\
\hline
\end{tabular}
}
  \caption{Amplifier size against profiling-set size. Identical loss, batch, epochs,
  split and seed in every cell. The zero-amplifier row is measured in Section 9.7; the
  four amplified cells are the ones discussed here.}
  \label{tab:ampdata}
\end{table}

The grid separates the two effects, and neither one dominates. More data helps at both
amplifier sizes, adding $0.222$ MSSIM at $72$ flops and $0.131$ at $4536$. A larger
amplifier helps at both data sizes, adding $0.227$ at $1.5$k traces and $0.136$ at $56$k.
They trade against each other along the diagonal: the $72$-flop amplifier with the paper's
profiling set reaches MSSIM $0.725$ and recognition $0.757$, which is close to what the
$63$ times larger amplifier reaches on the small profiling set, $0.730$ and $0.857$.

That is the correction to Section 9.5. Read on its own, the sweep says cutting the
amplifier $63$ times costs recognition $0.857 \to 0.400$ and makes the amplifier look
close to a precondition for the attack. Read with this grid, a $63$ times smaller
amplifier costs roughly what a $37$ times smaller profiling set costs, and either deficit
can be partly recovered by spending on the other. The amplifier changes how much
profiling data an attacker needs; it does not decide whether the attack is possible. For a
defender that is the worse reading of the two, because profiling traces are the cheap
resource in this threat model.

I should be careful not to overstate it in the other direction either, which I did briefly
when only the first two cells existed. The $72$-flop amplifier does not catch up with the
large one: at equal profiling budget the large amplifier is better in every column, and at
$56$k traces it reaches MSSIM $0.861$ and recognition $0.919$ against a ceiling of $0.898$
--- the recovered images classify slightly better than the originals, which is what
happens when a generator smooths its output toward class prototypes.

\subsubsection*{Against the published numbers}

Power2Picture report MSSIM $0.38$, pixel distance $77$ and recognition $0.53$ on
Fashion-MNIST. The corresponding numbers here are $0.861$, $14.5$ and $0.919$ at the large
amplifier, $0.725$, $23.7$ and $0.757$ at the smallest amplified build, and $0.724$,
$24.0$ and $0.746$ with the amplifier removed altogether (Section 9.7).

I would not present that as an improvement on their work, and it should not be read as
one. The training protocol now matches theirs, but the target does not: they attack a full
quantised LeNet through five on-chip TDC sensors on a shared FPGA, which is a remote attack
against a realistic accelerator, while this attacks a single convolution layer through an
external shunt with a leakage amplifier still present. Their measurement is harder in every
respect that matters. What the comparison does establish is that the pipeline in this
repository is not leaving anything obvious on the table relative to the published method,
which is what I wanted to know before drawing conclusions from it.

One practical note for anyone rerunning this. Loading $60\,000$ traces of $22\,144$ samples
needs $4.95$\,GB as float32, which does not fit alongside training; the loader now reduces
each shard to features as it reads it. I checked that this reproduces an earlier result
bit-identically before trusting it, having twice in this project been caught by a
refactor I assumed was safe.



\subsection*{9.7 The zero-amplifier point}

With the output registered, the bank-free design comes up, passes the functional check
bit-exact on all $676$ outputs, and captures normally. Running it under the same protocol
as the rest of the grid --- $60\,000$ Fashion training images, \texttt{gradmse}, batch
$256$, $20$ epochs, seed $0$, the same $56\,000/2000/2000$ split --- gives the row that was
missing.

The zero-amplifier row is the first line of Table \ref{tab:ampdata}.

Removing the amplifier entirely costs almost nothing. Against the $72$-flop build at the
same profiling set, MSSIM moves from $0.725$ to $0.724$ and recognition from $0.757$ to
$0.746$.

Those two runs are a paired comparison rather than two independent samples, which makes
the small differences easier to interpret. Both captures used the same shuffle seed and
the same start index, so they share a capture order; both trainings used the same split
seed. I checked that the recorded test indices are not merely the same size but element
for element identical, so both models were scored on the same $2000$ Fashion images in the
same order, and test-set sampling cannot account for the gap.

What remains is scoring noise and training stochasticity. For recognition the binomial
standard error on $2000$ images at $p \approx 0.75$ is $\sqrt{0.75 \cdot 0.25 / 2000}
\approx 0.0097$, so the observed $0.011$ difference is about one standard error --- I
would not read a direction into it. For MSSIM I make no noise claim and simply give the
comparison: removing the amplifier moves it by $0.001$, while going from $72$ to $4536$
flops in the same column moves it by $0.136$. The step that matters is two orders of
magnitude larger than the step to zero. That is the prediction the flattening implied, now measured rather than extrapolated,
and it removes the largest caveat this report carried: the attack does not need the
circuit I added. It needs profiling data.

I should be precise about what ``no amplifier'' means, because the honest number is not
zero flops. One flip-flop registers the leakage output, for the reason given above. It is
a single bit of state that no other logic reads, against $6624$ flops in the design as a
whole, and the utilisation reports separate the three builds cleanly: $6624$ flops here,
$6695$ with the $72$-flop bank, $11158$ with the $4536$-flop bank. The difference of $71$
is exactly the bank removed.

\subsubsection*{Checking it was really a different bitstream}

Mean trace amplitude is the usual discriminator between amplifier variants, and here it
almost fails: $2152.9$ against $2187.0$ for the $72$-flop build, a difference of $1.6$
percent. Given that an earlier sweep in this project was invalidated by four captures that
turned out to be one bitstream measured four times, a $1.6$ percent separation is not
evidence I would accept on its own. What establishes it instead is the SHA-256 the capture
manifest records for the programmed file, which differs from both the $72$-flop bitstream
and from the earlier dead build; the flop counts above; and the fact that the previous file
at that path read back all zeros, so a stale load would have aborted the run rather than
produced data. The small amplitude difference is itself consistent with Section 9.5 --- if
the curve has flattened by $216$ flops, the last $71$ should barely register.

\subsubsection*{What this does to the comparison with Power2Picture}

Table \ref{tab:papersettings} previously recorded the amplifier as a deviation from the
paper, listed as ``$72$ flops minimum''. That entry was wrong in a specific way worth
naming: the word ``minimum'' described a build of mine that would not synthesise into a
working design, and presented it as though it were a property of the hardware. With the
zero point measured the row now reads ``none'', and that deviation is closed.

Two deviations remain, and they are the substantial ones: the victim is a single
$3\times3$ convolution layer rather than a full quantised LeNet, and the measurement is an
external shunt rather than five on-chip TDC sensors. I have also added a third row to that
table which I had previously left out --- \texttt{DWELL}, which holds each convolution
window for eight cycles. A streaming accelerator advances one window per cycle, so that
parameter is an artifact of this bench in exactly the same way the amplifier was, and
leaving it unlisted while listing the amplifier was inconsistent.



\subsection*{9.8 Making the victim stream: \texttt{DWELL} = 2}

The amplifier was not the only thing in this bench that existed to make the attack easier.
\texttt{DWELL} holds each convolution window for that many clock cycles, and every result
above this point used $8$. A real streaming accelerator advances one window per cycle, so
a high dwell spreads the computation out and hands the attacker a longer, cleaner look at
each window. Having removed the amplifier because it was my addition, leaving \texttt{DWELL}
at $8$ without comment would have been inconsistent, so I rebuilt at $2$ --- the minimum
the core allows, since it requires an even dwell of at least two --- and captured again.

This is not a single-variable change, and it should not be read as one. At a fixed ADC
rate, holding each window for a quarter as long also gives a quarter as many samples of
it: $8$ per window instead of $32$, and a generator input of $1352$ features instead of
$5408$. Less integration time and a smaller input are the same physical change, and I
cannot separate them by turning one knob. Everything else is held: same $60\,000$-image
Fashion capture protocol, same $5$\,MHz clock and $40$\,dB gain, same loss, batch, epochs,
seed and split, no amplifier in either build, and --- checked rather than assumed --- the
same $2000$ test images in the same order, so this is a paired comparison.

\begin{table}[h]
  \centering
  \small
  \resizebox{\textwidth}{!}{\begin{tabular}{lrrrrrr}
\hline
Window dwell & Input dim. & MAE & Corr. & MSSIM & F1 excess & Recog. \\
\hline
$8$ cycles ($32$ samples/window) & 5408 & 23.98 & 0.878 & 0.724 & 0.632 & 0.746 \\
$2$ cycles ($8$ samples/window) & 1352 & 21.44 & 0.902 & 0.762 & 0.681 & 0.807 \\
\hline
\end{tabular}
}
  \caption{Window dwell at zero amplifier, $56$k training traces. Fewer cycles per window
  means fewer samples of it; both are consequences of the same change.}
  \label{tab:dwell}
\end{table}

I predicted this would get worse, and said so before running it. It got better, and by a
margin that is not noise: recognition rises from $0.746$ to $0.807$, which against the
binomial standard error of $0.0097$ on $2000$ images is a little over six standard errors,
and MSSIM rises from $0.724$ to $0.762$. Every metric in the table moves the same way. The
honest summary is that making the victim behave more like a real streaming accelerator did
not weaken this attack, it strengthened it.

\subsubsection*{Why, and what would settle it}

I have two candidate explanations and the data here does not choose between them.

The first is signal-to-noise. Compressing the same convolution into a quarter of the cycles
means more switching per cycle, and the traces show it directly: mean absolute amplitude is
$6546$ at \texttt{DWELL} $=2$ against $2153$ at \texttt{DWELL} $=8$, a factor of three, with
no amplifier in either. Fewer samples, each carrying a larger signal.

The second is the input dimension. The generator's first layer shrinks from $5408$ inputs
to $1352$, roughly a quarter of the parameters, while the training set stays at $56\,000$
images. A smaller input at a fixed data budget generalises better, and Section 9.6 already
established that this pipeline is data-limited rather than signal-limited in this regime.

These are cheap to separate and it needs no board. Re-featurising both existing captures in
\texttt{window\_abs} mode gives $676$ inputs for either dwell, equalising the dimension and
leaving amplitude as the only difference; if the $2$-cycle build still wins there, the
mechanism is signal-to-noise. I have not run it, and I am not going to assert a mechanism I
have not tested.

\subsubsection*{What this does to the comparison}

Two of the three deviations from Power2Picture that this report listed are now closed:
there is no leakage amplifier, and the victim no longer lingers on each window. Both were
mine, both were removed, and neither turned out to be holding the attack up --- the
amplifier cost essentially nothing to remove, and the dwell reduction helped. What remains
is the pair I cannot close cheaply: a single $3\times3$ convolution layer where they attack
a full quantised LeNet, and an external shunt where they read five on-chip TDC sensors. I
would expect the sensor change in particular to cost a great deal, since earlier
measurements in this project put that channel's correlation plateau near $0.33$.

For a defender the reading is the uncomfortable one. Two of the properties that looked most
like artifacts of a contrived test bench --- an added leakage circuit and an artificially
slow victim --- were the two that turned out not to matter.


\section*{10. Results}

\begin{table}[h]
  \centering
  \small
  \resizebox{\textwidth}{!}{\begin{tabular}{lrrrrrrrr}
\hline
Run & Pixel acc. & All-background & F1 & IoU & MSSIM & Pixel dist. & Recog. & $n$ \\
\hline
300 profile images, control evaluation set & 0.953 & 0.887 & 0.759 & 0.612 & 0.713 & 11.9 & 0.93 & 40 \\
40 profile / 40 evaluation, fresh images & 0.968 & 0.887 & 0.843 & 0.728 & 0.816 & 8.2 & 1.00 & 40 \\
40 profile / 40 evaluation, earlier run & 0.960 & 0.888 & 0.792 & 0.656 & 0.769 & 10.2 & 0.90 & 40 \\
Probe mismatch (one-hot template vs trained traces) & 0.811 & 0.888 & 0.119 & 0.063 & 0.094 & 48.1 & 0.20 & 40 \\
Single hardware repeat (no averaging) & 0.919 & 0.888 & 0.494 & 0.328 & 0.501 & 20.6 & 0.65 & 40 \\
\hline
\end{tabular}
}
  \caption{Main results. Every row is a real hardware capture. ``All-background'' is the
  trivial baseline and should always be compared with the pixel accuracy on its left.}
  \label{tab:headline}
\end{table}

\begin{table}[h]
  \centering
  \small
  \resizebox{\textwidth}{!}{\begin{tabular}{llrrccc}
\hline
Capture & Role & Clock & Gain & Not stock AES & Functional check & Mism. \\
\hline
rerun\_20260830\_trained\_80 & profiling (all condition tests) & 5.0\,MHz & 40\,dB & yes & 676/676 & 0 \\
rerun\_20260830\_onehot\_80 & profiling (probe mismatch) & 5.0\,MHz & 40\,dB & yes & 676/676 & 0 \\
prof\_5m & profiling (clock mixing, 5 MHz) & 5.0\,MHz & 40\,dB & yes & 676/676 & 0 \\
prof\_10m & profiling (clock mixing, 10 MHz) & 10.0\,MHz & 40\,dB & yes & 676/676 & 0 \\
grey\_p2p\_2000 & grey design, generative reconstruction & 5.0\,MHz & 40\,dB & yes & 676/676 & 0 \\
hard\_C1\_control & attack & 5.0\,MHz & 40\,dB & yes & 676/676 & 0 \\
hard\_C2\_reprogram & attack & 5.0\,MHz & 40\,dB & yes & 676/676 & 0 \\
hard\_C3\_clk7m5 & attack & 7.5\,MHz & 40\,dB & yes & 676/676 & 0 \\
hard\_C4\_clk10m & attack & 10.0\,MHz & 40\,dB & yes & 676/676 & 0 \\
hard\_C5\_gain30 & attack & 5.0\,MHz & 30\,dB & yes & 676/676 & 0 \\
hard\_C6\_gain50 & attack & 5.0\,MHz & 50\,dB & yes & 676/676 & 0 \\
paperscale\_500 & profiling only (300-image template) & 5.0\,MHz & 40\,dB & yes & 676/676 & 0 \\
\hline
\end{tabular}
}
  \caption{Hardware provenance of the captures used here. The functional check writes
  an image and a kernel to the FPGA, reads back the 676 outputs and compares them with a
  CPU model.}
\end{table}

\section*{11. What I understand from these numbers}

The attack works, and in the control condition it works well. The recovered images are
good enough that the digit recogniser reads them correctly.

Averaging is not a detail, it is the main practical cost of this attack on my board.
The convolution unit is small compared to the whole FPGA fabric, so its power is mixed
with the activity of everything else. This is the dilution problem. With one repeat the
attack is only slightly better than answering ``all background''. This also tells me
that the earlier low detection had two causes stacked on top of each other: the traces
came from the wrong design, and even correct traces need averaging.

The condition tests show where the method is fragile, and the fragility is specific:
the clock matters, the gain does not, and reprogramming the FPGA does not. This is
useful for us, because it is closer to a research contribution than repeating the
original paper. The original paper never reports what happens when the profiling
conditions and the attack conditions are different.

The paper-scale run taught me something I did not expect: more profiling data is not
automatically better. With a fixed search radius, a larger template returns more
candidates per window and the reconstruction becomes worse, even though the coverage of
the template improves. So the size of the template and the search radius have to be
tuned together. If I had only run the large template and reported it, I would have
concluded that my implementation is weaker than it really is.

That lesson turned out to be larger than the one run. The search radius is compared
against distances between normalised feature vectors, and those distances do not have the
same scale under different clocks, so a single radius quietly measures a different thing
in every condition. Once I tuned it per cell, the matched-condition result at 10\,MHz went
from $0.675$ to $0.882$ and the cost of clock randomisation went from $0.018$ to $0.145$.
Neither of those changed the direction of a conclusion, but both were large enough that
reporting them untuned was misleading, and I had done exactly that twice.

The result I understand least is also the one I find most interesting. I had assumed that
when a template fails under a changed clock, it fails because the power vectors no longer
match --- that is what ``the template is bound to its clock'' means, and it is what I wrote.
Section 8.7 says otherwise. In the pair where transfer fails, the true patch is retrieved
and ranked \emph{better} than in the pair where transfer works, and the reconstruction
still comes out near zero. So the template is not what breaks. Something in the greedy
assembly, which is the only stage that reasons about windows jointly rather than one at a
time, is throwing away an answer the earlier stages found. I cannot yet say what, but it
moves the interesting question away from the leakage and towards the algorithm, and it
means better features would not have helped.

\section*{12. What I suggest to do next}

Three items of the original list are done and are reported in Section 7. Two more that
appeared on this list in the previous version are now also done: the radius sweep is in
Section 8.1 and was extended to every cell of the grid in Section 8.3, and the
``profile at the higher clock'' rule was tested on four frequencies in Section 8.3, where
it turned out to be wrong as I had stated it. What is left, in the order I would do it:

\begin{enumerate}
\item Port the on-chip RO/TDC sensor into the grey design. With the amplifier gone, the
      measurement path is the largest remaining deviation from Power2Picture: they read
      five on-chip TDC sensors on a shared FPGA, a genuinely remote attack, while this
      reads an external shunt. The sensor RTL already exists in this repository but is
      wired into the BNN design rather than the grey one. Earlier measurements put the
      RO channel's correlation plateau near $0.33$, so this should be expected to make
      the reconstruction substantially worse, not better; that is the point, and it is
      the honest version of the threat model.
\item Find out what the greedy assembly stage is doing to the cross-condition cells.
      Section 8.7 is the most surprising result in this report: in the 5 and 10\,MHz pair,
      the direction that fails retrieves and ranks the true patch \emph{better} than the
      direction that works, so the loss is not in the template match. This is cheap to
      pursue because it needs no board --- the candidate lists can be dumped and the
      assembly replayed --- and it is the one place where I currently have a measured
      effect and no mechanism. If neighbouring windows really are disagreeing under a
      systematic offset, then a reconstruction that resolves disagreement globally instead
      of greedily should recover the upward cells, and that would be a genuine improvement
      on the paper's Algorithm 2 rather than a reproduction of it.
\item Recapture at 20\,MHz with two samples per cycle instead of four. Section 8.3 cannot
      separate weaker leakage at high clock from the ChipWhisperer-Lite's analog front end
      running at 80\,MSa/s, because both predict the same measurement. Halving the sample
      rate keeps the accelerator identical and relaxes the instrument, so if the
      correlation recovers the instrument was the limit. Until this is done I cannot say
      whether clocking the accelerator faster is a real countermeasure, which is the
      cheapest defence anyone would try.
\item Run the template attack against the grey design. Section 7.3 shows the grey
      levels are recoverable by the generative route, but the template attack still
      classifies binary patches and cannot express a grey value. Making it regress a
      grey patch instead is the missing piece for a direct comparison with the paper,
      which recovers grey levels with a template attack.
\item Redo the reprogramming test honestly, with a power cycle between profiling and
      attack, and put the row back in Table~2 with a real measurement behind it.
\item Measure transfer across days and temperatures, not only across settings inside one
      session.
\item Repeat the generative reconstruction on a capture made with averaging and nine
      probe kernels, so that it can be compared with the template attack fairly.
\end{enumerate}

One habit I want to carry forward, because it caught two wrong claims in this round. Both
the grid of Section 8.3 and the mixing sweep of Section 8.5 originally reported an optimum
that sat on the edge of the range I had swept, and in both cases I wrote a conclusion from
it without noticing. One of those conclusions was merely pessimistic; the other, in
Section 8.5, was a prediction from my own proposed mechanism that turned out to be wrong
when the range was extended. Checking that an optimum is bracketed before quoting it costs
nothing and would have caught both.

\section*{Appendix: how to reproduce}

All scripts are in the repository. The capture scripts require the board. The analysis
scripts do not.

\begin{quote}\small\ttfamily
host/cw305\_leakage\_capture.py \quad -- capture traces from the CW305\\
attack/hardtests\_20260830/run\_repeat\_sweep.sh \quad -- the averaging sweep\\
attack/hardtests\_20260830/run\_conditions.sh \quad -- the condition transfer tests\\
attack/hardtests\_20260830/run\_probe\_mismatch.sh \quad -- the probe mismatch test\\
attack/hardtests\_20260830/run\_paperscale.sh \quad -- captures 500 images and builds\\
\quad the 300-image template. Its own 200-image evaluation was abandoned; the run\\
\quad reported here attacks the 40 control images with that same template.\\
report/catchup\_v2\_2026-08-30/make\_figures.py \quad -- all figures\\
report/catchup\_v2\_2026-08-30/make\_tables.py \quad -- all tables
\end{quote}

Every figure that shows a power signal is produced from a real capture file. The tables
are generated directly from the score files, so no number in this report is typed by
hand.

Two notes about disk space, so that nothing looks missing. The machine ran out of space
during these runs. The \emph{attack bundles} are copies of the capture files with the
ground truth removed, so after each score was written I deleted the bundle; any of them
can be rebuilt with the \texttt{make-bundle} command from the capture directory. For the
paper-scale run I also deleted the bulk of the 500 raw captures after the template, the
bundle and the truth file had been extracted from them, and I kept the capture manifest
and three sample traces so that the provenance and the figures still work. Recapturing
those 500 images takes about five minutes on the bench.

\end{document}
